% ============================================================================
% MÓDULO 4. ANÁLISIS EXPLORATORIO Y ESTADÍSTICA PARA NEGOCIOS
% ============================================================================
\chapter{Análisis exploratorio y estadística para negocios}
\label{cap:eda}

\begin{objetivos}
\begin{itemize}
  \item Aplicar un proceso ordenado de \termino{análisis exploratorio de datos}
        (EDA) con una lista de verificación que puedes reutilizar en cualquier
        proyecto.
  \item Calcular e interpretar medidas de tendencia central, dispersión,
        percentiles y forma de una distribución, y saber cuándo el promedio
        engaña.
  \item Detectar datos atípicos con reglas objetivas y decidir qué hacer con
        ellos sin borrarlos a ciegas.
  \item Construir KPIs de desempeño (crecimiento, acumulados, promedios
        móviles, índices y avance contra meta) y un análisis de Pareto/ABC.
  \item Medir relaciones con correlaciones y distinguir correlación de
        causalidad.
  \item Usar las distribuciones normal, binomial y de Poisson para responder
        preguntas de probabilidad del día a día.
  \item Estimar con intervalos de confianza, aplicar pruebas de hipótesis
        (t, proporciones, ANOVA, chi-cuadrada) y ajustar regresiones lineales,
        explicando los resultados en lenguaje de negocio.
\end{itemize}
\end{objetivos}

Imagina que trabajas como analista en una cadena de tiendas con diez
sucursales en México. Es lunes y la directora comercial te escribe:
\emph{``Antes de aprobar el presupuesto de fin de año necesito entender qué
está pasando con las ventas: ¿cuánto gasta un cliente típico?, ¿hay tiendas
que funcionen mejor que otras?, ¿nuestros clientes premium realmente gastan
más?, ¿qué productos sostienen el negocio? Y si invertimos más en publicidad,
¿cuánto deberíamos esperar vender?''}.

Con lo que aprendiste en los Módulos~\ref{cap:manipulacion}
y~\ref{cap:visualizacion} ya puedes filtrar, agrupar, unir tablas y graficar.
Pero la directora no quiere tablas sueltas: quiere \emph{respuestas}. Para
darlas necesitas dos cosas. La primera es un \textbf{método} para recorrer los
datos sin perderte: eso es el análisis exploratorio. La segunda es la
\textbf{estadística}, que te permite distinguir un patrón real de una simple
casualidad y decir con honestidad cuánta confianza tienes en cada número.

En este módulo aprenderás ambas cosas con los datos del curso. No necesitas
saber matemáticas avanzadas: cada fórmula aparece explicada en palabras, se
calcula en R paso a paso y termina con la pregunta que más importa en una
empresa: \emph{¿y esto qué significa para el negocio?}

% ----------------------------------------------------------------------------
\section{¿Qué es el análisis exploratorio (EDA)?}
\label{sec:m4-eda}

\begin{concepto}{Análisis exploratorio de datos (EDA)}
El \termino{análisis exploratorio de datos} (EDA, por sus siglas en inglés,
\emph{Exploratory Data Analysis}) es el proceso de conocer un conjunto de
datos \emph{antes} de sacar conclusiones: revisar su estructura, su calidad,
cómo se reparten los valores, cómo se relacionan las variables y qué
preguntas vale la pena investigar. El término lo popularizó el estadístico
John Tukey en 1977 con una idea sencilla: primero mira los datos, después
decide qué prueba aplicar.
\end{concepto}

Piensa en el EDA como la consulta de un médico: antes de recetar, toma la
presión, pregunta síntomas y pide análisis. Si te saltas ese paso, puedes
recetar el medicamento equivocado. En BI pasa lo mismo: un dashboard
construido sin EDA puede mostrar con mucha elegancia un número que está mal.

La figura~\ref{fig:m4-proceso} resume el proceso que seguiremos en todo el
módulo. No es una receta rígida: a menudo un hallazgo en el paso 5 te obliga a
volver al paso 2.

\begin{figure}[H]
\centering
\begin{tikzpicture}[
  paso/.style={draw=primario, fill=primario!6, rounded corners=3pt,
               minimum width=2.3cm, minimum height=1cm, align=center,
               font=\sffamily\footnotesize},
  flecha/.style={-{Stealth[length=2.5mm]}, thick, color=gristexto}]
  \node[paso] (p1) {1. Estructura};
  \node[paso, right=0.45cm of p1] (p2) {2. Calidad};
  \node[paso, right=0.45cm of p2] (p3) {3. Univariado};
  \node[paso, right=0.45cm of p3] (p4) {4. Bivariado};
  \node[paso, below=0.9cm of p4] (p5) {5. Segmentos};
  \node[paso, left=0.45cm of p5] (p6) {6. Atípicos};
  \node[paso, left=0.45cm of p6, fill=acento!10, draw=acento] (p7)
       {7. Hipótesis};
  \node[paso, left=0.45cm of p7, fill=exito!10, draw=exito] (p8)
       {Pregunta de\\negocio};
  \draw[flecha] (p1) -- (p2);
  \draw[flecha] (p2) -- (p3);
  \draw[flecha] (p3) -- (p4);
  \draw[flecha] (p4) -- (p5);
  \draw[flecha] (p5) -- (p6);
  \draw[flecha] (p6) -- (p7);
  \draw[flecha] (p7) -- (p8);
  \draw[flecha, dashed] (p8) -- (p1);
\end{tikzpicture}
\caption{Proceso de análisis exploratorio. La flecha punteada indica que cada
hipótesis genera nuevas preguntas y el ciclo se repite.}
\label{fig:m4-proceso}
\end{figure}

La tabla~\ref{tab:m4-checklist} convierte el proceso en una lista de
verificación. Te recomiendo copiarla en tus proyectos: evita que se te olvide
revisar algo importante.

\begin{table}[H]
\centering
\small
\caption{Lista de verificación (\emph{checklist}) del análisis exploratorio.}
\label{tab:m4-checklist}
\begin{tabularx}{\textwidth}{@{}l X l@{}}
\toprule
\textbf{Paso} & \textbf{Pregunta que respondes} & \textbf{Funciones de R} \\
\midrule
1. Estructura & ¿Cuántas filas y columnas hay? ¿Qué tipo tiene cada
  columna? ¿Qué representa cada fila? & \codigo{dim()}, \codigo{glimpse()} \\
2. Calidad & ¿Hay valores faltantes, duplicados, negativos imposibles o
  categorías mal escritas? & \codigo{skim()}, \codigo{summary()},
  \codigo{count()} \\
3. Univariado & ¿Cómo se reparte cada variable? ¿Cuál es su valor típico y
  cuánto varía? & \codigo{mean()}, \codigo{quantile()},
  \codigo{geom\_histogram()} \\
4. Bivariado & ¿Qué variables se mueven juntas? & \codigo{cor()},
  \codigo{geom\_point()} \\
5. Segmentos & ¿Cambian los resultados por tienda, categoría o tipo de
  cliente? & \codigo{group\_by()} + \codigo{summarise()} \\
6. Atípicos & ¿Hay valores extremos? ¿Son errores o casos valiosos? &
  \codigo{IQR()}, \codigo{geom\_boxplot()} \\
7. Hipótesis & ¿Qué diferencias parecen reales y cuáles pueden ser azar? &
  \codigo{t.test()}, \codigo{lm()} \\
\bottomrule
\end{tabularx}
\end{table}

\subsection{Preparar el entorno}

En este módulo usaremos dos paquetes nuevos: \paquete{skimr}, que resume un
data frame completo con una sola función, y \paquete{corrplot}, que dibuja
matrices de correlación. Instálalos una sola vez:

\begin{rnoejecutar}
install.packages(c("skimr", "corrplot"))
\end{rnoejecutar}

Ahora carga los paquetes y los datos. Recuerda que el directorio de trabajo
debe ser la carpeta del curso (la que contiene \archivo{datasets/}).

\begin{rcode}
# Paquetes del módulo
library(tidyverse)   # dplyr, ggplot2, readr, tidyr... (Módulos 2 y 3)
library(skimr)       # resúmenes rápidos de un data frame
library(corrplot)    # matrices de correlación
library(scales)      # formatos de negocio: dollar(), percent(), comma()

# Datos del curso
transacciones <- read_csv("datasets/transacciones.csv",
                          show_col_types = FALSE)
clientes  <- read_csv("datasets/clientes.csv", show_col_types = FALSE)
productos <- read_csv("datasets/productos.csv", show_col_types = FALSE)
tiendas   <- read_csv("datasets/tiendas.csv", show_col_types = FALSE)
\end{rcode}

También definimos los colores corporativos y el tema gráfico que construiste
en el Módulo~\ref{cap:visualizacion}, para que todas las gráficas del módulo
tengan el mismo estilo:

\begin{rcode}
# Colores corporativos
color_principal <- "#2a78d6"   # una sola serie
color_resalte   <- "#eb6834"   # para destacar un elemento
color_gris      <- "#b5b3ad"   # el resto, cuando se resalta uno
paleta_libro <- c("#2a78d6", "#eb6834", "#1baf7a", "#eda100",
                  "#e87ba4", "#008300", "#4a3aa7", "#e34948")

# Tema corporativo (misma forma que en el Módulo 3)
tema_libro <- function(base_size = 11) {
  theme_minimal(base_size = base_size) +
    theme(
      plot.title       = element_text(face = "bold", color = "#0b0b0b"),
      plot.subtitle    = element_text(color = "#52514e"),
      plot.caption     = element_text(color = "#52514e", size = rel(0.8)),
      axis.title       = element_text(color = "#52514e"),
      axis.text        = element_text(color = "#52514e"),
      panel.grid.minor = element_blank(),
      panel.grid.major = element_line(color = "#e6e5e1", linewidth = 0.3),
      legend.position  = "top",
      legend.title     = element_text(color = "#52514e")
    )
}
\end{rcode}

\subsection{Pasos 1 y 2: estructura y calidad}

Empieza siempre por lo básico: ¿qué tamaño tiene la tabla y qué contiene cada
columna?

\begin{rcode}
dim(transacciones)        # filas y columnas
glimpse(transacciones)    # tipo y primeros valores de cada columna
\end{rcode}
\begin{rsalida}
[1] 1000   14
Rows: 1,000
Columns: 14
$ transaccion_id    <dbl> 1, 2, 3, 4, 5, 6, 7, 8, 9, 10, 11, 12, 13,~
$ fecha             <date> 2023-01-13, 2023-01-09, 2023-02-20, 2023-~
$ hora              <time> 13:15:00, 10:45:00, 10:30:00, 10:45:00, 1~
$ cliente_id        <dbl> 155, 155, 75, 51, 173, 135, 90, 168, 13, 9~
$ producto_id       <dbl> 3, 47, 6, 17, 28, 46, 30, 48, 21, 12, 49, ~
$ cantidad          <dbl> 4, 4, 3, 4, 1, 2, 4, 2, 2, 5, 2, 3, 4, 2, ~
$ precio_venta      <dbl> 303.97, 58.07, 275.46, 105.28, 205.47, 534~
$ metodo_pago       <chr> "Tarjeta Crédito", "Tarjeta Débito", "Tarj~
$ tienda_id         <dbl> 10, 2, 7, 6, 7, 7, 8, 8, 5, 2, 7, 6, 8, 2,~
$ vendedor_id       <dbl> 4, 22, 14, 5, 14, 22, 23, 5, 3, 14, 5, 9, ~
$ total_transaccion <dbl> 1215.88, 232.28, 826.38, 421.12, 205.47, 1~
$ mes               <chr> "2023-01", "2023-01", "2023-02", "2023-05"~
$ trimestre         <chr> "Q1", "Q1", "Q1", "Q2", "Q1", "Q1", "Q2", ~
$ dia_semana        <chr> "viernes", "lunes", "lunes", "viernes", "m~
\end{rsalida}

Cada fila es una \emph{transacción} (una línea de venta): un cliente compró
cierta \codigo{cantidad} de un producto a un \codigo{precio\_venta} en una
tienda. La columna \codigo{total\_transaccion} es cantidad $\times$ precio.
Fíjate en que \codigo{read\_csv()} reconoció \codigo{fecha} como fecha
(\codigo{<date>}) y \codigo{hora} como hora (\codigo{<time>}); los
identificadores (\codigo{cliente\_id}, \codigo{tienda\_id}\ldots) son
números, pero en realidad son \emph{etiquetas}: promediarlos no tiene
sentido.

La función \codigo{summary()} da un primer resumen numérico. Aplícala solo a
las columnas de medida para que la salida sea legible:

\begin{rcode}
transacciones %>%
  select(cantidad, precio_venta, total_transaccion) %>%
  summary()
\end{rcode}
\begin{rsalida}
    cantidad      precio_venta    total_transaccion
 Min.   :1.000   Min.   : 20.14   Min.   :  22.65
 1st Qu.:2.000   1st Qu.:164.11   1st Qu.: 341.11
 Median :3.000   Median :324.92   Median : 713.49
 Mean   :2.969   Mean   :318.31   Mean   : 959.99
 3rd Qu.:4.000   3rd Qu.:468.00   3rd Qu.:1478.62
 Max.   :5.000   Max.   :598.81   Max.   :2992.35
\end{rsalida}

Para cada columna obtienes el mínimo (\codigo{Min.}), el primer cuartil
(\codigo{1st Qu.}), la mediana (\codigo{Median}), la media (\codigo{Mean}), el
tercer cuartil (\codigo{3rd Qu.}) y el máximo (\codigo{Max.}). Ya aparece una
pista importante: en \codigo{total\_transaccion} la media (959.99) es bastante
mayor que la mediana (713.49). En la sección~\ref{sec:m4-forma} verás que eso
indica una distribución con ``cola'' hacia los valores altos.

El paquete \paquete{skimr} va más allá: con \codigo{skim()} obtienes en una
sola llamada el número de faltantes (\codigo{n\_missing}), el porcentaje de
datos completos (\codigo{complete\_rate}), la media, la desviación estándar y
los percentiles de cada columna, agrupados por tipo de dato. En la consola de
RStudio, \codigo{skim()} dibuja además un mini histograma con caracteres
especiales; para un reporte impreso como este libro usamos su variante
\codigo{skim\_without\_charts()}, que muestra exactamente lo mismo sin ese
dibujo.

\begin{rcode}
skim_without_charts(transacciones)
\end{rcode}
\begin{rsalida}
-- Data Summary ------------------------
                           Values
Name                       transacciones
Number of rows             1000
Number of columns          14
_______________________
Column type frequency:
  character                4
  Date                     1
  difftime                 1
  numeric                  8
________________________
Group variables            None

-- Variable type: character ------------------------------------------
  skim_variable n_missing complete_rate min max empty n_unique
1 metodo_pago           0             1   8  15     0        4
2 mes                   0             1   7   7     0       12
3 trimestre             0             1   2   2     0        4
4 dia_semana            0             1   5   9     0        7
  whitespace
1          0
2          0
3          0
4          0

-- Variable type: Date -----------------------------------------------
  skim_variable n_missing complete_rate min        max
1 fecha                 0             1 2023-01-01 2023-12-31
  median     n_unique
1 2023-06-23      345

-- Variable type: difftime -------------------------------------------
  skim_variable n_missing complete_rate min        max
1 hora                  0             1 28800 secs 74700 secs
  median     n_unique
1 51300 secs       52

-- Variable type: numeric --------------------------------------------
  skim_variable     n_missing complete_rate   mean     sd   p0  p25
1 transaccion_id            0             1 500.   289.    1   251.
2 cliente_id                0             1  97.8   57.5   1    49
3 producto_id               0             1  25.3   14.4   1    13
4 cantidad                  0             1   2.97   1.45  1     2
5 precio_venta              0             1 318.   172.   20.1 164.
6 tienda_id                 0             1   5.61   2.88  1     3
7 vendedor_id               0             1  13.2    7.21  1     7
8 total_transaccion         0             1 960.   760.   22.6 341.
   p50    p75  p100
1 500.  750.  1000
2  97   150    200
3  25    37.2   50
4   3     4      5
5 325.  468.   599.
6   6     8     10
7  13    19     25
8 713. 1479.  2992.
\end{rsalida}

\begin{ejemplo}{Leer un \texttt{skim()} como analista}
¿Qué te dice esta salida en 30 segundos?
\begin{itemize}
  \item \textbf{Calidad}: \codigo{n\_missing} es 0 en todas las columnas y
        \codigo{complete\_rate} es 1: no hay faltantes. En datos reales casi
        nunca pasa; cuando veas un \codigo{complete\_rate} de 0.85, sabes que
        15\% de esa columna está vacía.
  \item \textbf{Categorías}: \codigo{metodo\_pago} tiene 4 valores distintos
        (\codigo{n\_unique}) y \codigo{mes}, 12. Si \codigo{metodo\_pago}
        tuviera 7 valores, sospecharías de errores de captura como
        ``tarjeta credito'' y ``Tarjeta Crédito''.
  \item \textbf{Cobertura}: las fechas van del 1 de enero al 31 de diciembre de
        2023, pero solo hay 345 fechas distintas: hubo días sin ninguna venta
        registrada.
  \item \textbf{Rangos}: \codigo{cantidad} va de 1 a 5 y el ticket
        (\codigo{total\_transaccion}) de 22.6 a 2,992; no hay negativos ni
        ceros imposibles.
\end{itemize}
\end{ejemplo}

\begin{hazlotu}
Aplica \codigo{skim\_without\_charts()} a \codigo{clientes} y a
\codigo{productos}. ¿Cuántos segmentos de edad hay? ¿Hay algún producto con
\codigo{margen\_pct} negativo? (Pista: mira el valor \codigo{p0}, que es el
mínimo.)
\end{hazlotu}

% ----------------------------------------------------------------------------
\section{Estadística descriptiva}
\label{sec:m4-descriptiva}

La \termino{estadística descriptiva} resume muchos datos en pocos números.
Responde dos preguntas: ¿cuál es el valor \emph{típico}? (tendencia central)
y ¿qué tanto se alejan los datos de ese valor típico? (dispersión). Para no
repetir \codigo{transacciones\$total\_transaccion} en cada línea, guardamos
esa columna en un vector llamado \codigo{ticket}:

\begin{rcode}
ticket <- transacciones$total_transaccion   # 1,000 tickets de venta
length(ticket)
\end{rcode}
\begin{rsalida}
[1] 1000
\end{rsalida}

\subsection{Medidas de tendencia central: media, mediana y moda}

La \termino{media} (o promedio) suma todos los valores y divide entre cuántos
son. Si llamamos $x_1, x_2, \ldots, x_n$ a los $n$ tickets:
\[
  \bar{x} \;=\; \frac{x_1 + x_2 + \cdots + x_n}{n}
          \;=\; \frac{1}{n}\sum_{i=1}^{n} x_i
\]
El símbolo $\sum$ (sigma) solo significa ``suma todo''. La \termino{mediana}
es el valor que queda justo a la mitad cuando ordenas los datos de menor a
mayor: la mitad de los tickets está por debajo y la otra mitad por encima. La
\termino{moda} es el valor que más se repite; es útil sobre todo para
variables categóricas (método de pago, día de la semana).

\begin{rcode}
media_ticket   <- mean(ticket)     # promedio
mediana_ticket <- median(ticket)   # valor central
media_ticket
mediana_ticket
sum(ticket) / length(ticket)       # la media "a mano": mismo resultado
\end{rcode}
\begin{rsalida}
[1] 959.9927
[1] 713.485
[1] 959.9927
\end{rsalida}

R no trae una función para la moda (la función \codigo{mode()} de R hace otra
cosa: dice el tipo de un objeto). La construimos nosotros con lo que
aprendiste en el Módulo~\ref{cap:fundamentos}:

\begin{rcode}
# Moda: el valor que más veces aparece en un vector
moda <- function(x) {
  frecuencias <- table(x)                      # cuenta cada valor
  names(frecuencias)[which.max(frecuencias)]   # el más frecuente
}

moda(transacciones$metodo_pago)   # método de pago más usado
moda(transacciones$dia_semana)    # día con más transacciones
moda(transacciones$cantidad)      # unidades más comunes por ticket
\end{rcode}
\begin{rsalida}
[1] "Tarjeta Crédito"
[1] "viernes"
[1] "1"
\end{rsalida}

\codigo{table()} cuenta cuántas veces aparece cada valor y
\codigo{which.max()} devuelve la posición del mayor conteo. La moda siempre
sale como texto (por eso \codigo{"1"} aparece entre comillas); si la
necesitas como número usa \codigo{as.numeric(moda(x))}. Si hay empate,
\codigo{which.max()} devuelve el primero en orden alfabético, así que
conviene revisar la tabla completa:

\begin{rcode}
table(transacciones$cantidad)
\end{rcode}
\begin{rsalida}
  1   2   3   4   5
219 196 190 187 208
\end{rsalida}

\begin{negocio}{La moda en la operación}
La moda responde preguntas operativas muy concretas: ¿cuál es la talla que más
se vende? (para surtir más), ¿cuál es el método de pago más usado? (para
negociar comisiones con el banco), ¿a qué hora llegan más clientes? (para
programar turnos). Aquí, la tarjeta de crédito es el método más frecuente y
lo más común es que un ticket lleve 1 sola unidad.
\end{negocio}

\subsection{Medidas de dispersión}

Dos tiendas pueden tener el mismo ticket promedio y comportarse de forma muy
distinta: en una todos gastan alrededor de 900 pesos y en la otra unos gastan
50 y otros 3,000. La \termino{dispersión} mide esa diferencia.

\begin{itemize}
  \item \termino{Rango}: máximo menos mínimo. Fácil, pero depende solo de dos
        datos (los más extremos).
  \item \termino{Varianza} ($s^2$): el promedio de las distancias al
        cuadrado respecto a la media. Elevamos al cuadrado para que las
        distancias negativas no cancelen a las positivas:
        \[ s^2 = \frac{1}{n-1}\sum_{i=1}^{n}(x_i - \bar{x})^2 \]
        Se divide entre $n-1$ y no entre $n$ por un ajuste técnico que
        corrige un pequeño sesgo cuando trabajas con muestras; con muchos
        datos la diferencia es mínima.
  \item \termino{Desviación estándar} ($s$): la raíz cuadrada de la
        varianza. Regresa a las unidades originales (pesos), por eso es la
        medida de dispersión más usada. Se interpreta como ``cuánto se aleja,
        típicamente, un dato del promedio''.
  \item \termino{Rango intercuartílico} (IQR): distancia entre el percentil
        75 y el 25. Mide la amplitud del 50\% central de los datos y no le
        afectan los extremos.
  \item \termino{Coeficiente de variación} (CV): la desviación estándar
        dividida entre la media, $CV = s/\bar{x}$. Como no tiene unidades,
        sirve para comparar la variabilidad de cosas distintas (tickets en
        pesos contra unidades vendidas).
\end{itemize}

\begin{rcode}
range(ticket)                  # mínimo y máximo
diff(range(ticket))            # rango = máximo - mínimo
var(ticket)                    # varianza (pesos al cuadrado)
sd(ticket)                     # desviación estándar (pesos)
IQR(ticket)                    # rango intercuartílico
sd(ticket) / mean(ticket)      # coeficiente de variación

# La varianza "a mano", para que veas que no hay magia:
sum((ticket - mean(ticket))^2) / (length(ticket) - 1)
\end{rcode}
\begin{rsalida}
[1]   22.65 2992.35
[1] 2969.7
[1] 578152.5
[1] 760.3634
[1] 1137.513
[1] 0.7920512
[1] 578152.5
\end{rsalida}

¿Qué nos dice esto? Un ticket típico se aleja unos \$760 del promedio de
\$960. El coeficiente de variación es 0.79: la desviación estándar equivale al
79\% de la media, lo que indica tickets \emph{muy} heterogéneos. Como regla
práctica, un CV menor a 0.15 indica datos bastante homogéneos, entre 0.15 y
0.30 variabilidad moderada y por encima de 0.30 alta variabilidad. La
varianza (578,152.5) está en ``pesos al cuadrado'', una unidad sin sentido de
negocio; por eso casi siempre reportarás la desviación estándar.

\begin{hazlotu}
Calcula el coeficiente de variación de \codigo{cantidad} y de
\codigo{precio\_venta}. ¿Cuál de las dos varía más en términos relativos?
\end{hazlotu}

\subsection{Percentiles y cuantiles}

El \termino{percentil} $p$ es el valor que deja por debajo al $p\%$ de los
datos. La mediana es el percentil 50; el primer y tercer
\termino{cuartil} son los percentiles 25 y 75. En R se calculan con
\codigo{quantile()}, indicando las proporciones que quieres:

\begin{rcode}
quantile(ticket, probs = c(0.10, 0.25, 0.50, 0.75, 0.90, 0.95))
\end{rcode}
\begin{rsalida}
      10%       25%       50%       75%       90%       95%
 159.1200  341.1125  713.4850 1478.6250 2113.6140 2501.2725
\end{rsalida}

Los percentiles se leen así: el 10\% de los tickets es menor a \$159.12; la
mitad, menor a \$713.49; y solo el 5\% supera \$2,501.27. Son perfectos para
fijar umbrales de negocio: si quieres un programa de beneficios para el 10\%
de compras más grandes, el corte está en \$2,113.61 (percentil 90).

\begin{ejemplo}{Umbral para un programa de lealtad}
La gerencia de mercadotecnia quiere regalar un cupón a los tickets que estén
en el 20\% más alto. ¿A partir de qué monto se entrega el cupón y cuántos
tickets del año lo habrían recibido?
\end{ejemplo}

\begin{rcode}
umbral_cupon <- quantile(ticket, probs = 0.80)   # percentil 80
umbral_cupon
sum(ticket > umbral_cupon)   # tickets que superan el umbral
\end{rcode}
\begin{rsalida}
     80%
1641.174
[1] 200
\end{rsalida}

El cupón se entregaría en compras mayores a \$1,641.17 y, con los datos de
2023, lo habrían recibido 200 tickets (el 20\% de 1,000, como era de
esperarse). Si cada cupón cuesta 50 pesos, el costo del programa sería de
unos 10,000 pesos al año.

\subsection{Cuándo la media engaña}
\label{sec:m4-media-engana}

La media tiene un defecto: \textbf{es muy sensible a los valores extremos}.
Un solo dato gigante la arrastra hacia arriba; la mediana, en cambio, apenas
se mueve. Veámoslo con dos ejemplos pequeños creados con \codigo{tibble()}.

\begin{ejemplo}{Un distribuidor con un cliente gigante}
Un distribuidor mayorista tiene 10 clientes. Nueve compran entre 20 y 30 mil
pesos al mes; el décimo es una cadena nacional que compra 900 mil. El
director pregunta: ``¿cuánto compra un cliente típico?''.
\end{ejemplo}

\begin{rcode}
compras_b2b <- tibble(
  cliente = paste("Cliente", 1:10),
  compra_mensual = c(22000, 25000, 21000, 28000, 30000,
                     24000, 26000, 23000, 27000, 900000)
)

compras_b2b %>%
  summarise(media   = mean(compra_mensual),
            mediana = median(compra_mensual))
\end{rcode}
\begin{rsalida}
# A tibble: 1 x 2
   media mediana
   <dbl>   <dbl>
1 112600   25500
\end{rsalida}

Si respondes ``un cliente típico compra 112,600 pesos'' estarás describiendo
a un cliente que no existe: nueve de diez compran menos de 31 mil. La mediana
(25,500) sí describe al cliente típico. La media sigue siendo útil para otra
pregunta: si multiplicas 112,600 por 10 clientes, obtienes exactamente las
ventas totales. \textbf{La media sirve para totales; la mediana, para
describir al caso típico.}

El segundo ejemplo es clásico en recursos humanos:

\begin{rcode}
salarios <- c(rep(15000, 9), 250000)   # 9 empleados y un director
mean(salarios)
median(salarios)
\end{rcode}
\begin{rsalida}
[1] 38500
[1] 15000
\end{rsalida}

Anunciar que ``el salario promedio en la empresa es de 38,500 pesos'' es
técnicamente cierto, pero engañoso: nueve de cada diez personas ganan 15,000.
Por eso las estadísticas oficiales de ingresos suelen reportar la mediana.

\begin{consejo}
Reporta siempre la media \emph{y} la mediana. Si son parecidas, los datos son
bastante simétricos y cualquiera sirve. Si la media es claramente mayor que la
mediana, hay valores altos que la arrastran: usa la mediana para describir el
caso típico y la media solo para calcular totales.
\end{consejo}

\subsection{Resumen estadístico por grupo}

En el Módulo~\ref{cap:manipulacion} aprendiste \codigo{group\_by()} +
\codigo{summarise()}. Ahora lo combinamos con todas las medidas anteriores
para comparar segmentos. Pregunta de negocio: \emph{¿el comportamiento del
ticket cambia según el método de pago?}

\begin{rcode}
resumen_pago <- transacciones %>%
  group_by(metodo_pago) %>%
  summarise(
    n        = n(),                                 # cuántos tickets
    ventas   = round(sum(total_transaccion)),       # total vendido
    media    = mean(total_transaccion),             # ticket promedio
    mediana  = median(total_transaccion),           # ticket típico
    desv_est = sd(total_transaccion),               # dispersión
    cv       = sd(total_transaccion) / mean(total_transaccion),
    .groups = "drop"
  ) %>%
  mutate(pct_ventas = ventas / sum(ventas)) %>%     # % de las ventas
  arrange(desc(ventas))

resumen_pago
\end{rcode}
\begin{rsalida}
# A tibble: 4 x 8
  metodo_pago         n ventas media mediana desv_est    cv pct_ventas
  <chr>           <int>  <dbl> <dbl>   <dbl>    <dbl> <dbl>      <dbl>
1 Tarjeta Crédito   350 351171 1003.    780.     782. 0.779     0.366
2 Efectivo          315 303735  964.    710.     757. 0.785     0.316
3 Tarjeta Débito    233 220855  948.    732.     738. 0.779     0.230
4 Transferencia     102  84232  826.    608.     741. 0.897     0.0877
\end{rsalida}

Lectura de negocio: la tarjeta de crédito concentra la mayor parte de las
ventas (36.6\%) y tiene el ticket promedio más alto (\$1,003); la
transferencia es el método menos usado y el de ticket más bajo (\$826). Los
cuatro métodos tienen un CV alto (de 0.78 a 0.90): en todos conviven tickets
muy pequeños y muy grandes. ¿La diferencia entre \$1,003 y \$826 es real o es azar? No lo
sabemos todavía: en la sección~\ref{sec:m4-hipotesis} aprenderás a
responderlo con una prueba de hipótesis.

\begin{hazlotu}
Repite el resumen agrupando por \codigo{tienda\_id}. ¿Qué tienda tiene la
mediana de ticket más alta? ¿Coincide con la que tiene la media más alta?
\end{hazlotu}

% ----------------------------------------------------------------------------
\section{Forma de la distribución}
\label{sec:m4-forma}

Los números resumen ayudan, pero nada sustituye a \emph{ver} los datos. La
\termino{distribución} de una variable describe qué valores toma y con qué
frecuencia. Su herramienta básica es el \termino{histograma}: divide el eje en
intervalos y cuenta cuántos datos caen en cada uno. Encima podemos dibujar una
\termino{curva de densidad}, una versión suavizada del histograma que deja
ver la forma general sin el ``ruido'' de las barras.

\begin{ejemplo}{¿Cómo se reparte el ticket de compra?}
La directora quiere saber si ``el ticket típico'' es un buen resumen o si hay
grupos de clientes con compras muy distintas. Dibujamos el histograma del
ticket con su curva de densidad y marcamos la media y la mediana.
\end{ejemplo}

\begin{rcode}
p_hist <- ggplot(transacciones, aes(x = total_transaccion)) +
  # barras: cuántos tickets caen en cada intervalo de 150 pesos
  geom_histogram(binwidth = 150, boundary = 0,
                 fill = color_principal, color = "white", alpha = 0.85) +
  # densidad reescalada (x 150) para que quede en la escala de las barras
  geom_density(aes(y = after_stat(count) * 150),
               color = "#0b0b0b", linewidth = 0.8) +
  # media (línea continua) y mediana (línea punteada)
  geom_vline(xintercept = media_ticket, color = color_resalte,
             linewidth = 0.8) +
  geom_vline(xintercept = mediana_ticket, color = color_resalte,
             linewidth = 0.8, linetype = "dashed") +
  annotate("text", x = media_ticket + 50, y = 140, hjust = 0,
           label = paste("Media:", dollar(media_ticket)),
           color = color_resalte, size = 3.5) +
  annotate("text", x = mediana_ticket - 50, y = 140, hjust = 1,
           label = paste("Mediana:", dollar(mediana_ticket)),
           color = color_resalte, size = 3.5) +
  scale_x_continuous(labels = dollar) +
  labs(title = "Distribución del ticket por transacción",
       subtitle = "Muchos tickets pequeños y una cola larga a la derecha",
       x = "Total de la transacción", y = "Número de transacciones") +
  tema_libro()

p_hist
\end{rcode}

\begin{figure}[H]
\centering
\includegraphics[width=0.9\textwidth]{m04-histograma-ticket}
\caption{Histograma y curva de densidad del ticket. La media queda a la
derecha de la mediana porque los tickets grandes la ``jalan''.}
\label{fig:m04-histograma-ticket}
\end{figure}

La figura~\ref{fig:m04-histograma-ticket} confirma lo que sugerían los
números: la mayoría de los tickets se concentra por debajo de 600 pesos y, a
partir de ahí, las barras bajan poco a poco formando una \termino{cola}
larga hacia la derecha. A esa forma se le llama \termino{asimetría
positiva} (o sesgo a la derecha). Observa que la media queda a la derecha de
la mediana: los pocos tickets grandes la arrastran.

\begin{consejo}
El ancho de las barras (\codigo{binwidth}) cambia mucho la impresión que da un
histograma. Prueba siempre dos o tres anchos distintos antes de sacar
conclusiones; con barras muy anchas se esconden patrones y con barras muy
angostas todo parece ruido.
\end{consejo}

\subsection{Medir la asimetría}

Podemos resumir la forma con un número: el \termino{coeficiente de
asimetría}. La idea es promediar las distancias a la media \emph{elevadas al
cubo}. Al elevar al cubo se conserva el signo (las distancias negativas
siguen siendo negativas) y los valores lejanos pesan muchísimo; si la cola
derecha es más larga, dominan los cubos positivos. Para que no dependa de las
unidades, se divide entre $s^3$:
\[
  \text{asimetría} \;=\;
  \frac{\frac{1}{n}\sum_{i=1}^{n}(x_i - \bar{x})^3}{s^3}
\]

\begin{rcode}
# Coeficiente de asimetría (versión sencilla)
asimetria <- function(x) {
  mean((x - mean(x))^3) / sd(x)^3
}

asimetria(ticket)                        # ticket de compra
asimetria(transacciones$precio_venta)    # precio unitario
asimetria(compras_b2b$compra_mensual)    # distribuidor con cliente gigante
\end{rcode}
\begin{rsalida}
[1] 0.8711898
[1] -0.08471191
[1] 2.276413
\end{rsalida}

La tabla~\ref{tab:m4-asimetria} te ayuda a interpretarlo. El ticket tiene una
asimetría de 0.87 (positiva y moderada); el precio unitario está muy cerca de
cero (los precios se reparten de forma pareja entre 20 y 600 pesos); y el
ejemplo del distribuidor, con un cliente gigante, tiene una asimetría fuerte
de 2.28.

\begin{table}[H]
\centering
\small
\caption{Cómo interpretar el coeficiente de asimetría.}
\label{tab:m4-asimetria}
\begin{tabularx}{\textwidth}{@{}l X X@{}}
\toprule
\textbf{Valor} & \textbf{Forma} & \textbf{Ejemplo de negocio} \\
\midrule
Entre $-0.5$ y $0.5$ & Aproximadamente simétrica; media $\approx$ mediana &
  Estatura de clientes, tiempos de entrega estables \\
Entre $0.5$ y $1$ & Asimetría positiva moderada & Ticket de compra de
  este módulo \\
Mayor a $1$ & Asimetría positiva fuerte (cola derecha larga) & Ingresos de
  clientes, ventas por producto en catálogos grandes \\
Menor a $-0.5$ & Asimetría negativa (cola izquierda) & Calificaciones de
  satisfacción cuando casi todos ponen 9 o 10 \\
\bottomrule
\end{tabularx}
\end{table}

\begin{negocio}{El ticket con cola larga}
Una distribución con cola larga a la derecha es la norma en el comercio: hay
muchas compras chicas y pocas muy grandes. Esto tiene tres consecuencias
prácticas:
\begin{enumerate}
  \item \textbf{Metas}: si fijas la meta del ticket con la media, la mayoría
        de los vendedores quedarán por debajo aunque trabajen bien; la mediana
        es una referencia más justa para el ``vendedor típico''.
  \item \textbf{Estrategia}: la cola derecha son los clientes que más valen.
        Identificarlos (percentil 90 o 95) permite diseñar atención
        preferente.
  \item \textbf{Oportunidad}: el gran bloque de tickets bajos es el terreno
        para estrategias de venta cruzada (``¿le agrego\ldots?'') que suban el
        ticket típico.
\end{enumerate}
\end{negocio}

\begin{hazlotu}
Dibuja el histograma de \codigo{precio\_venta} con \codigo{binwidth = 50} y
calcula su asimetría. ¿Se parece a la forma del ticket? ¿Por qué crees que el
ticket (cantidad $\times$ precio) tiene cola y el precio no?
\end{hazlotu}

% ----------------------------------------------------------------------------
\section{Datos atípicos (outliers)}
\label{sec:m4-atipicos}

Un \termino{dato atípico} (en inglés, \emph{outlier}) es un valor que se
aleja mucho del resto. Puede ser un \textbf{error} (alguien tecleó 29,923.50
en vez de 2,992.35), un \textbf{caso legítimo pero excepcional} (una empresa
que compró 50 laptops) o una \textbf{señal de fraude}. El trabajo del analista
no es borrarlo: es \emph{detectarlo, investigarlo y decidir}.

\subsection{Regla de 1.5 $\times$ IQR}

La regla más usada la propuso también John Tukey y es la que dibujan los
diagramas de caja (\emph{boxplots}). Se consideran atípicos los valores que
caen fuera de estas ``cercas'':
\[
  \text{límite inferior} = Q_1 - 1.5 \times IQR
  \qquad\qquad
  \text{límite superior} = Q_3 + 1.5 \times IQR
\]
donde $Q_1$ y $Q_3$ son los percentiles 25 y 75, e $IQR = Q_3 - Q_1$. Como
usa percentiles, la regla no se deja engañar por los propios atípicos.

\begin{rcode}
q1  <- quantile(ticket, 0.25)
q3  <- quantile(ticket, 0.75)
iqr <- q3 - q1
limite_inferior <- unname(q1 - 1.5 * iqr)   # unname() quita el "25%"
limite_superior <- unname(q3 + 1.5 * iqr)
c(limite_inferior, limite_superior)

# ¿Cuántos tickets quedan fuera de las cercas?
sum(ticket < limite_inferior | ticket > limite_superior)
\end{rcode}
\begin{rsalida}
[1] -1365.156  3184.894
[1] 0
\end{rsalida}

Ningún ticket es atípico según esta regla: el límite superior (\$3,184.89)
está por encima del ticket máximo (\$2,992.35). Tiene sentido: una transacción
tiene como máximo 5 unidades de un producto de hasta 600 pesos. El límite
inferior es negativo, lo cual simplemente significa que por abajo no hay
forma de ser atípico. \textbf{``No hay atípicos'' también es un resultado
valioso}: significa que puedes usar los tickets con tranquilidad.

\subsection{Puntuación z}

La segunda regla usa la \termino{puntuación z} (\emph{z-score}): cuántas
desviaciones estándar se aleja cada dato de la media.
\[
  z_i = \frac{x_i - \bar{x}}{s}
\]
Un $z = 2$ significa ``dos desviaciones estándar por encima del promedio''.
Se suele marcar como atípico todo dato con $|z| > 3$.

\begin{rcode}
z_ticket <- (ticket - mean(ticket)) / sd(ticket)
range(z_ticket)              # el z más bajo y el más alto
sum(abs(z_ticket) > 3)       # tickets con |z| > 3
\end{rcode}
\begin{rsalida}
[1] -1.232756  2.672876
[1] 0
\end{rsalida}

El ticket más alto está a 2.67 desviaciones estándar de la media: tampoco hay
atípicos por esta regla. Para reutilizar la regla del IQR la convertimos en
función:

\begin{rcode}
# Devuelve TRUE para los valores fuera de las cercas de Tukey
marcar_atipicos <- function(x) {
  q1 <- quantile(x, 0.25, na.rm = TRUE)
  q3 <- quantile(x, 0.75, na.rm = TRUE)
  rango_iq <- q3 - q1
  x < q1 - 1.5 * rango_iq | x > q3 + 1.5 * rango_iq
}
\end{rcode}

\subsection{Cuando los atípicos son errores}

Para ver las reglas en acción con errores de verdad, \textbf{simulamos} una
situación frecuente: al exportar del punto de venta se colaron tres registros
problemáticos. Agregamos a una copia de los tickets una venta con el punto
decimal recorrido (29,923.50), una devolución capturada como venta negativa
($-450$) y un ticket en cero.

\begin{rcode}
# Copia de los tickets + 3 registros con problemas (simulados)
tickets_revision <- transacciones %>%
  select(transaccion_id, total_transaccion) %>%
  bind_rows(tibble(transaccion_id    = c(1001, 1002, 1003),
                   total_transaccion = c(29923.50, -450, 0)))

tickets_revision %>%
  mutate(
    atipico_iqr = marcar_atipicos(total_transaccion),
    z = (total_transaccion - mean(total_transaccion)) /
          sd(total_transaccion),
    atipico_z = abs(z) > 3
  ) %>%
  filter(atipico_iqr | atipico_z | total_transaccion <= 0)
\end{rcode}
\begin{rsalida}
# A tibble: 3 x 5
  transaccion_id total_transaccion atipico_iqr      z atipico_z
           <dbl>             <dbl> <lgl>        <dbl> <lgl>
1           1001            29924. TRUE        24.3   TRUE
2           1002             -450  FALSE       -1.21  FALSE
3           1003                0  FALSE       -0.829 FALSE
\end{rsalida}

Tres lecciones importantes:
\begin{enumerate}
  \item Ambas reglas detectan el error de captura (29,923.50): su $z$ es
        24.3, absurdamente lejos.
  \item \textbf{Ninguna regla estadística detecta el $-450$ ni el 0}, porque
        están cerca del centro de los datos. Para eso necesitas
        \termino{reglas de negocio}: ``un ticket de venta debe ser mayor a
        cero''. La estadística no sustituye al conocimiento del negocio.
  \item Un solo error contamina las medidas sensibles:
\end{enumerate}

\begin{rcode}
sd(transacciones$total_transaccion)       # sin los registros erróneos
sd(tickets_revision$total_transaccion)    # con los registros erróneos
\end{rcode}
\begin{rsalida}
[1] 760.3634
[1] 1189.887
\end{rsalida}

La desviación estándar pasó de 760 a 1,190 (un 56\% más), sobre todo por
\emph{un} solo dato erróneo entre 1,003. Por eso la puntuación $z$ es menos confiable que la regla del
IQR cuando hay atípicos fuertes: el propio atípico infla la $s$ con la que se
le mide.

\subsection{Cuando los atípicos son tus mejores clientes}

Ahora un caso real. Pregunta de negocio: \emph{¿hay clientes con un gasto
anual fuera de lo normal?} Calculamos el gasto de cada cliente en 2023 y
aplicamos las dos reglas.

\begin{rcode}
gasto_clientes <- transacciones %>%
  group_by(cliente_id) %>%
  summarise(compras = n(),
            gasto   = sum(total_transaccion),
            .groups = "drop") %>%
  mutate(atipico = marcar_atipicos(gasto),
         z       = (gasto - mean(gasto)) / sd(gasto))

gasto_clientes %>%
  filter(atipico) %>%
  arrange(desc(gasto))
\end{rcode}
\begin{rsalida}
# A tibble: 8 x 5
  cliente_id compras  gasto atipico     z
       <dbl>   <int>  <dbl> <lgl>   <dbl>
1          5      13 14680. TRUE     3.65
2         90       9 12504. TRUE     2.85
3        183       9 11848. TRUE     2.60
4        168      12 11676. TRUE     2.54
5         55      13 11356. TRUE     2.42
6        125      10 11240. TRUE     2.38
7        162      10 11176. TRUE     2.35
8         31       7 10980. TRUE     2.28
\end{rsalida}

La regla del IQR marca 8 clientes con gasto anual superior a unos 10,800
pesos; la regla de $|z| > 3$ solo marcaría al cliente 5. Ninguno parece un
error: todos compraron entre 7 y 13 veces en el año. \textbf{Son tus mejores
clientes}, no basura estadística. Borrarlos del análisis sería un error grave:
son justo los que un programa de lealtad debe cuidar.

\begin{rcode}
set.seed(123)   # geom_jitter() mueve los puntos al azar
p_caja <- ggplot(gasto_clientes, aes(x = gasto, y = "")) +
  geom_boxplot(width = 0.5, outlier.shape = NA, fill = "grey96") +
  geom_jitter(aes(color = atipico), height = 0.15,
              size = 1.8, alpha = 0.8) +
  scale_color_manual(values = c(`FALSE` = color_gris,
                                `TRUE` = color_resalte),
                     labels = c("Dentro de las cercas",
                                "Atípico (regla 1.5 × IQR)"),
                     name = NULL) +
  scale_x_continuous(labels = dollar) +
  labs(title = "Gasto anual por cliente en 2023",
       subtitle = "8 clientes superan el límite superior de la caja",
       x = "Gasto total del cliente", y = NULL) +
  tema_libro()

p_caja
\end{rcode}

\begin{figure}[H]
\centering
\includegraphics[width=0.9\textwidth]{m04-boxplot-clientes}
\caption{Diagrama de caja del gasto anual por cliente. Cada punto es un
cliente; en naranja, los que la regla del IQR marca como atípicos.}
\label{fig:m04-boxplot-clientes}
\end{figure}

En la figura~\ref{fig:m04-boxplot-clientes} la caja va del percentil 25 al 75,
la línea interior es la mediana y los ``bigotes'' llegan hasta el último dato
dentro de las cercas. Usamos \codigo{outlier.shape = NA} para que
\codigo{geom\_boxplot()} no dibuje sus propios puntos atípicos, porque ya los
dibujamos todos con \codigo{geom\_jitter()}.

\subsection{¿Qué hacer con un atípico?}

\begin{table}[H]
\centering
\small
\caption{Qué hacer con un dato atípico según su origen.}
\label{tab:m4-atipicos}
\begin{tabularx}{\textwidth}{@{}l X X@{}}
\toprule
\textbf{Situación} & \textbf{Acción} & \textbf{Ejemplo} \\
\midrule
Error de captura evidente & Corregir en la fuente si es posible; si no,
  excluir y documentarlo & 29,923.50 en lugar de 2,992.35 \\
Viola una regla de negocio & Separar y reportar al área responsable &
  Tickets negativos o en cero \\
Caso legítimo excepcional & Conservar; analizar aparte o usar medidas
  robustas (mediana, IQR) & Compra corporativa de 50 laptops \\
Cliente o producto de alto valor & Conservar y estudiar: suele ser una
  oportunidad & Los 8 clientes del ejemplo anterior \\
Posible fraude & Escalar a auditoría antes de tocarlo & Muchas devoluciones
  del mismo vendedor \\
\bottomrule
\end{tabularx}
\end{table}

\begin{advertencia}
Nunca borres atípicos ``porque estorban'' para que la gráfica se vea bonita
o el promedio salga mejor. Si decides excluir datos, deja constancia en tu
script (un comentario y un \codigo{filter()} explícito) y menciónalo en el
reporte: ``se excluyeron 3 registros con errores de captura''.
\end{advertencia}

\begin{hazlotu}
Carga \archivo{ventas\_retail.csv} y aplica \codigo{marcar\_atipicos()} a la
columna \codigo{total}. ¿Cuántas ventas quedan fuera de las cercas? Míralas
con \codigo{filter()}: ¿son errores o ventas grandes legítimas?
\end{hazlotu}

% ----------------------------------------------------------------------------
\section{KPIs y métricas de desempeño}
\label{sec:m4-kpis}

Un \termino{KPI} (\emph{Key Performance Indicator}, indicador clave de
desempeño) es una métrica que la empresa sigue de forma periódica para saber
si va bien o mal. Las ventas del mes son un dato; las ventas del mes
\emph{comparadas} contra el mes anterior, contra la meta o contra el mismo
mes del año pasado son un KPI. La tabla~\ref{tab:m4-kpis} resume los cálculos
más comunes.

\begin{table}[H]
\centering
\small
\caption{KPIs de desempeño más usados y su cálculo en R.}
\label{tab:m4-kpis}
\begin{tabularx}{\textwidth}{@{}l X l@{}}
\toprule
\textbf{KPI} & \textbf{Qué responde} & \textbf{Cálculo en R} \\
\midrule
Crecimiento MoM & ¿Cuánto crecimos contra el mes anterior?
  (\emph{Month over Month}) & \codigo{ventas / lag(ventas) - 1} \\
Acumulado (YTD) & ¿Cuánto llevamos vendido en el año?
  (\emph{Year to Date}) & \codigo{cumsum(ventas)} \\
Promedio móvil & ¿Cuál es la tendencia sin el ruido de cada mes? &
  \codigo{(v + lag(v) + lag(v, 2)) / 3} \\
Índice base 100 & ¿Cuánto cambiamos respecto a un punto de partida? &
  \codigo{ventas / first(ventas) * 100} \\
Variación vs meta & ¿Estamos arriba o abajo de lo planeado? &
  \codigo{ventas / meta - 1} \\
\bottomrule
\end{tabularx}
\end{table}

\begin{ejemplo}{Tablero mensual de ventas}
La meta comercial de 2023 era vender \$80,000 al mes. Construye una tabla
con las ventas mensuales, su crecimiento contra el mes anterior, el
acumulado del año, el promedio móvil de 3 meses, el índice con enero = 100 y
la variación contra la meta.
\end{ejemplo}

\begin{rcode}
meta_mensual <- 80000   # meta de ventas por mes

kpi_mensual <- transacciones %>%
  group_by(mes) %>%
  summarise(ventas = sum(total_transaccion), .groups = "drop") %>%
  arrange(mes) %>%                                  # orden cronológico
  mutate(
    crec_mom   = ventas / lag(ventas) - 1,          # vs mes anterior
    acumulado  = cumsum(ventas),                    # acumulado del año
    prom_movil = (ventas + lag(ventas) + lag(ventas, 2)) / 3,
    indice_100 = ventas / first(ventas) * 100,      # enero = 100
    vs_meta    = ventas / meta_mensual - 1          # sobre/bajo la meta
  )

kpi_mensual
\end{rcode}
\begin{rsalida}
# A tibble: 12 x 7
   mes     ventas crec_mom acumulado prom_movil indice_100  vs_meta
   <chr>    <dbl>    <dbl>     <dbl>      <dbl>      <dbl>    <dbl>
 1 2023-01 82560.  NA         82560.        NA       100    0.0320
 2 2023-02 65899.  -0.202    148459         NA        79.8 -0.176
 3 2023-03 90931.   0.380    239390.     79797.      110.   0.137
 4 2023-04 57038.  -0.373    296427.     71289.       69.1 -0.287
 5 2023-05 93910.   0.646    390338.     80626.      114.   0.174
 6 2023-06 88509.  -0.0575   478847.     79819.      107.   0.106
 7 2023-07 86456.  -0.0232   565303.     89625.      105.   0.0807
 8 2023-08 78929.  -0.0871   644232.     84632.       95.6 -0.0134
 9 2023-09 68856.  -0.128    713088.     78080.       83.4 -0.139
10 2023-10 70004.   0.0167   783092.     72596.       84.8 -0.125
11 2023-11 80507.   0.150    863599.     73122.       97.5  0.00634
12 2023-12 96394.   0.197    959993.     82302.      117.   0.205
\end{rsalida}

Cómo se lee cada columna, tomando marzo (\codigo{2023-03}) como ejemplo:
\begin{itemize}
  \item \codigo{crec\_mom} = 0.380: marzo vendió 38\% más que febrero.
  \item \codigo{acumulado} = 239,390: entre enero y marzo se vendieron
        \$239,390.
  \item \codigo{prom\_movil} = 79,797: el promedio de enero, febrero y marzo.
        Los dos primeros meses quedan en \codigo{NA} porque todavía no hay
        tres meses para promediar. La función \codigo{lag()} ``trae'' el
        valor de la fila anterior, y \codigo{lag(ventas, 2)} el de dos filas
        antes.
  \item \codigo{indice\_100} = 110: marzo vendió 10\% más que enero.
  \item \codigo{vs\_meta} = 0.137: marzo superó la meta en 13.7\%.
\end{itemize}

Para presentarla a la dirección conviene darle formato de negocio con las
funciones de \paquete{scales}:

\begin{rcode}
kpi_mensual %>%
  transmute(
    mes,
    ventas     = dollar(ventas, accuracy = 1),
    crec_mom   = percent(crec_mom, accuracy = 0.1),
    acumulado  = dollar(acumulado, accuracy = 1),
    vs_meta    = percent(vs_meta, accuracy = 0.1)
  )
\end{rcode}
\begin{rsalida}
# A tibble: 12 x 5
   mes     ventas  crec_mom acumulado vs_meta
   <chr>   <chr>   <chr>    <chr>     <chr>
 1 2023-01 $82,560 <NA>     $82,560   3.2%
 2 2023-02 $65,899 -20.2%   $148,459  -17.6%
 3 2023-03 $90,931 38.0%    $239,390  13.7%
 4 2023-04 $57,038 -37.3%   $296,427  -28.7%
 5 2023-05 $93,910 64.6%    $390,338  17.4%
 6 2023-06 $88,509 -5.8%    $478,847  10.6%
 7 2023-07 $86,456 -2.3%    $565,303  8.1%
 8 2023-08 $78,929 -8.7%    $644,232  -1.3%
 9 2023-09 $68,856 -12.8%   $713,088  -13.9%
10 2023-10 $70,004 1.7%     $783,092  -12.5%
11 2023-11 $80,507 15.0%    $863,599  0.6%
12 2023-12 $96,394 19.7%    $959,993  20.5%
\end{rsalida}

¿Qué nos dice el tablero? Los resultados mes a mes brincan mucho ($-20.2\%$,
$+38.0\%$, $-37.3\%$, $+64.6\%$\ldots). Abril (\$57,038) fue el peor mes,
28.7\% por debajo de la meta, y diciembre el mejor, 20.5\% por encima.
Contemos cuántos meses cumplieron la meta y cuánto faltó o sobró en el año:

\begin{rcode}
meses_arriba <- sum(kpi_mensual$vs_meta >= 0)   # meses que cumplieron
meses_arriba
sum(kpi_mensual$ventas) - meta_mensual * 12     # diferencia anual
\end{rcode}
\begin{rsalida}
[1] 7
[1] -7.31
\end{rsalida}

La meta mensual se cumplió en 7 de los 12 meses y el año cerró con \$959,993:
apenas \$7.31 por debajo de la meta anual de \$960,000. En el total anual la
empresa está en la meta, pero con un camino muy irregular. Para ver la
tendencia sin tanto ruido graficamos las ventas junto con su promedio móvil.

\begin{rcode}
# Formato largo: una fila por mes y por serie (Módulo 2, pivot_longer)
kpi_largo <- kpi_mensual %>%
  mutate(fecha_mes = as.Date(paste0(mes, "-01"))) %>%
  select(fecha_mes,
         `Ventas del mes` = ventas,
         `Promedio móvil 3 meses` = prom_movil) %>%
  pivot_longer(-fecha_mes, names_to = "serie", values_to = "valor")

p_kpi <- ggplot(kpi_largo, aes(x = fecha_mes, y = valor, color = serie)) +
  geom_hline(yintercept = meta_mensual, linetype = "dashed",
             color = "#52514e", linewidth = 0.5) +
  annotate("text", x = as.Date("2023-01-01"), y = meta_mensual + 2500,
           label = "Meta: $80,000", hjust = 0, size = 3.3,
           color = "#52514e") +
  geom_line(linewidth = 0.8, na.rm = TRUE) +
  geom_point(size = 1.8, na.rm = TRUE) +
  scale_color_manual(
    values = c("Ventas del mes" = color_principal,
               "Promedio móvil 3 meses" = color_resalte),
    breaks = c("Ventas del mes", "Promedio móvil 3 meses"),
    name = NULL) +
  scale_y_continuous(labels = dollar) +
  scale_x_date(date_breaks = "1 month", date_labels = "%b") +
  labs(title = "Ventas mensuales 2023 contra la meta",
       subtitle = "El promedio móvil suaviza los brincos de cada mes",
       x = NULL, y = "Ventas") +
  tema_libro()

p_kpi
\end{rcode}

\begin{figure}[H]
\centering
\includegraphics[width=0.9\textwidth]{m04-kpi-evolucion}
\caption{Ventas mensuales, promedio móvil de 3 meses y meta mensual.}
\label{fig:m04-kpi-evolucion}
\end{figure}

La figura~\ref{fig:m04-kpi-evolucion} cuenta la historia completa: las ventas
mensuales (azul) suben y bajan con fuerza, mientras que el promedio móvil
(naranja) muestra la tendencia de fondo: rondó la meta en primavera (con el
bache de abril), la superó en julio y agosto, cayó por debajo durante el
otoño y se recuperó en diciembre.

\begin{negocio}{Por qué los directivos piden promedios móviles}
Un mes malo aislado puede deberse a días festivos, a un problema de
inventario o simplemente al azar. Reaccionar a cada brinco (``abril cayó
37\%: ¡recorten gastos!'') lleva a decisiones erráticas. El promedio móvil
filtra ese ruido: si \emph{también} baja durante varios meses, entonces sí
hay una tendencia que atender. En empresas con estacionalidad fuerte se
complementa con el crecimiento contra el mismo mes del año anterior
(\emph{YoY}, \emph{Year over Year}).
\end{negocio}

\begin{consejo}
Para promedios móviles de muchos periodos (por ejemplo, 12 meses), escribir
\codigo{lag()} doce veces es incómodo. El paquete \paquete{zoo} (que se
instala junto con \paquete{forecast}) tiene
\codigo{zoo::rollmean(ventas, k = 12, fill = NA, align = "right")}, que hace
exactamente el mismo cálculo.
\end{consejo}

\begin{hazlotu}
Calcula los mismos KPIs por \codigo{trimestre} (crecimiento contra el
trimestre anterior, acumulado e índice base 100). ¿Qué trimestre fue el
mejor? ¿Cuál fue el crecimiento del cuarto trimestre contra el tercero?
\end{hazlotu}

% ----------------------------------------------------------------------------
\section{Análisis de Pareto y clasificación ABC}
\label{sec:m4-pareto}

El \termino{principio de Pareto}, o regla 80/20, dice que en muchos negocios
una minoría de elementos genera la mayoría de los resultados: el 20\% de los
productos genera el 80\% de las ventas, el 20\% de los clientes aporta el
80\% de los ingresos. Los porcentajes exactos varían; lo importante es la
idea de \emph{concentración}. El \termino{análisis ABC} aprovecha esa idea
para clasificar productos:

\begin{itemize}
  \item \textbf{Clase A}: los productos que, sumados de mayor a menor,
        acumulan el primer 80\% de los ingresos. Son vitales.
  \item \textbf{Clase B}: los que acumulan del 80\% al 95\%. Importantes.
  \item \textbf{Clase C}: el 5\% restante. Muchos productos que aportan poco.
\end{itemize}

\begin{ejemplo}{¿Qué productos sostienen el negocio?}
La directora quiere saber si hay un grupo pequeño de productos ``estrella''
del que dependa la empresa, para asegurarse de que nunca falten en
inventario.
\end{ejemplo}

\begin{rcode}
pareto <- transacciones %>%
  group_by(producto_id) %>%
  summarise(ingresos = sum(total_transaccion), .groups = "drop") %>%
  arrange(desc(ingresos)) %>%                        # de mayor a menor
  mutate(
    pct_ingresos  = ingresos / sum(ingresos),        # peso de cada uno
    pct_acumulado = cumsum(pct_ingresos),            # % acumulado
    pct_productos = row_number() / n(),              # % del catálogo
    clase = case_when(
      pct_acumulado <= 0.80 ~ "A",
      pct_acumulado <= 0.95 ~ "B",
      TRUE                  ~ "C"
    )
  ) %>%
  left_join(productos %>%
              select(producto_id, categoria, estado_stock, activo),
            by = "producto_id")

pareto %>%
  select(producto_id, ingresos, pct_ingresos, pct_acumulado, clase) %>%
  head(10)
\end{rcode}
\begin{rsalida}
# A tibble: 10 x 5
   producto_id ingresos pct_ingresos pct_acumulado clase
         <dbl>    <dbl>        <dbl>         <dbl> <chr>
 1           6   34028.       0.0354        0.0354 A
 2          19   28368.       0.0296        0.0650 A
 3          31   25430.       0.0265        0.0915 A
 4          47   24527.       0.0255        0.117  A
 5          46   24470.       0.0255        0.143  A
 6          32   24333.       0.0253        0.168  A
 7          14   24023.       0.0250        0.193  A
 8          24   23593.       0.0246        0.217  A
 9          29   22809.       0.0238        0.241  A
10          21   22083.       0.0230        0.264  A
\end{rsalida}

El producto más vendido (el 6) aporta apenas el 3.5\% de los ingresos, y los
10 primeros productos (el 20\% del catálogo) suman el 26.4\%. Muy lejos de
un 80/20. Veamos el resumen por clase:

\begin{rcode}
pareto %>%
  group_by(clase) %>%
  summarise(productos     = n(),
            pct_catalogo  = n() / nrow(pareto),
            ingresos      = sum(ingresos),
            pct_ingresos  = sum(pct_ingresos),
            .groups = "drop")
\end{rcode}
\begin{rsalida}
# A tibble: 3 x 5
  clase productos pct_catalogo ingresos pct_ingresos
  <chr>     <int>        <dbl>    <dbl>        <dbl>
1 A            36         0.72  764155.       0.796
2 B             9         0.18  137845.       0.144
3 C             5         0.1    57992.       0.0604
\end{rsalida}

\begin{rcode}
p_pareto <- ggplot(pareto, aes(x = pct_productos, y = pct_acumulado)) +
  # diagonal: lo que pasaría si todos los productos vendieran lo mismo
  geom_abline(slope = 1, intercept = 0, linetype = "dashed",
              color = color_gris) +
  geom_hline(yintercept = 0.80, color = "#52514e", linewidth = 0.4) +
  geom_line(color = "#52514e", linewidth = 0.6) +
  geom_point(aes(color = clase), size = 2) +
  # punto de referencia de la regla 80/20
  annotate("point", x = 0.2, y = 0.8, shape = 4, size = 4, stroke = 1.2) +
  annotate("text", x = 0.22, y = 0.86, hjust = 0, size = 3.3,
           label = "Regla 80/20 (referencia)") +
  scale_color_manual(values = paleta_libro[1:3], name = "Clase") +
  scale_x_continuous(labels = percent) +
  scale_y_continuous(labels = percent, limits = c(0, 1)) +
  labs(title = "Curva de Pareto de ingresos por producto",
       subtitle = "Ingresos muy repartidos: no hay pocos productos estrella",
       x = "% de productos (de mayor a menor ingreso)",
       y = "% acumulado de ingresos") +
  tema_libro()

p_pareto
\end{rcode}

\begin{figure}[H]
\centering
\includegraphics[width=0.9\textwidth]{m04-pareto}
\caption{Curva de Pareto. Cuanto más se despega la curva de la diagonal
punteada, más concentrados están los ingresos.}
\label{fig:m04-pareto}
\end{figure}

\begin{ejemplo}{Interpretar la curva de Pareto}
En la figura~\ref{fig:m04-pareto}, el eje horizontal es el porcentaje del
catálogo (ordenado de mayor a menor ingreso) y el vertical, el porcentaje
acumulado de ingresos. Si la empresa siguiera la regla 80/20, la curva
pasaría por la cruz. Nuestra curva está muy pegada a la diagonal: los 36
productos de la clase A (el 72\% del catálogo) apenas juntan el 79.6\% de los
ingresos. \textbf{Este es un hallazgo en sí mismo}: la empresa no depende de
unos cuantos productos estrella; sus ingresos están muy repartidos.
\end{ejemplo}

¿Qué implica para el inventario? Cruzamos la clase con el estado del stock:

\begin{rcode}
pareto %>% count(clase, estado_stock)

# Productos A con stock bajo: prioridad de reabastecimiento
pareto %>%
  filter(clase == "A", estado_stock != "Stock OK") %>%
  select(producto_id, categoria, ingresos, estado_stock)
\end{rcode}
\begin{rsalida}
# A tibble: 4 x 3
  clase estado_stock     n
  <chr> <chr>        <int>
1 A     Stock Bajo       4
2 A     Stock OK        32
3 B     Stock OK         9
4 C     Stock OK         5
# A tibble: 4 x 4
  producto_id categoria  ingresos estado_stock
        <dbl> <chr>         <dbl> <chr>
1           6 Juguetes     34028. Stock Bajo
2          33 Alimentos    21240. Stock Bajo
3           3 Automotriz   19908. Stock Bajo
4          39 Alimentos    18337. Stock Bajo
\end{rsalida}

\begin{negocio}{Políticas de inventario ABC}
\begin{itemize}
  \item \textbf{Clase A} (36 productos): revisión semanal de inventario,
        stock de seguridad alto y nunca permitir faltantes. Los 4 productos A
        con stock bajo (6, 33, 3 y 39) deben reabastecerse de inmediato;
        entre ellos está el producto 6, ¡el más vendido del año!
  \item \textbf{Clase B} (9 productos): revisión quincenal y stock de
        seguridad moderado.
  \item \textbf{Clase C} (5 productos): stock mínimo, pedidos bajo demanda y
        candidatos a revisión (¿se descontinúan?, ¿se promueven?).
\end{itemize}
Como nuestro catálogo es tan parejo, sacar productos C ahorraría poco espacio
pero tampoco perdería muchas ventas: 5 productos que suman el 6.0\% de los
ingresos. En catálogos reales, con miles de productos, la clase C suele
tener cientos de artículos que ocupan bodega y casi no se venden: ahí el ABC
libera mucho dinero.
\end{negocio}

\begin{hazlotu}
Haz el mismo análisis de Pareto con los \emph{clientes} en lugar de los
productos. ¿Qué porcentaje de clientes genera el 80\% de los ingresos? ¿Los
ingresos están más o menos concentrados que en los productos?
\end{hazlotu}

% ----------------------------------------------------------------------------
\section{Correlación: ¿qué variables se mueven juntas?}
\label{sec:m4-correlacion}

La \termino{correlación} mide qué tanto se mueven juntas dos variables
numéricas. Se resume en un número entre $-1$ y $+1$:
\begin{itemize}
  \item Cerca de $+1$: cuando una sube, la otra también sube (visitas a la
        tienda y ventas).
  \item Cerca de $-1$: cuando una sube, la otra baja (precio y unidades
        vendidas).
  \item Cerca de $0$: no hay relación \emph{lineal} (estatura del cliente y
        ticket).
\end{itemize}
Como guía práctica: $|r| < 0.3$ es una relación débil, entre $0.3$ y $0.7$
moderada, y mayor a $0.7$ fuerte.

\subsection{Pearson y Spearman}

La \termino{correlación de Pearson} ($r$) mide relaciones \emph{lineales}
(en línea recta) y usa los valores tal cual; por eso es sensible a los
atípicos. Su fórmula compara cómo se desvía cada variable de su media:
\[
  r = \frac{\sum (x_i - \bar{x})(y_i - \bar{y})}
           {\sqrt{\sum (x_i - \bar{x})^2}\,\sqrt{\sum (y_i - \bar{y})^2}}
\]
Si cuando $x$ está por encima de su media $y$ también lo está, los productos
del numerador son positivos y $r$ sale positiva. La \termino{correlación de
Spearman} ($\rho$, ``ro'') hace lo mismo pero con los \emph{lugares}
(rangos) que ocupa cada dato al ordenarlos; mide si la relación es
\emph{monótona} (siempre sube o siempre baja, aunque no sea en línea recta) y
casi no le afectan los atípicos.

\begin{ejemplo}{¿Qué explica el gasto de un cliente?}
Construimos una tabla con una fila por cliente: número de compras, unidades,
gasto total, ticket promedio, edad y antigüedad como cliente. Queremos saber
qué variables se relacionan con el gasto anual.
\end{ejemplo}

\begin{rcode}
perfil_clientes <- transacciones %>%
  group_by(cliente_id) %>%
  summarise(compras   = n(),
            unidades  = sum(cantidad),
            gasto     = sum(total_transaccion),
            ticket    = mean(total_transaccion),
            .groups = "drop") %>%
  left_join(clientes, by = "cliente_id") %>%
  mutate(antiguedad = as.numeric(as.Date("2023-12-31") - fecha_registro))

# Pearson (el método por defecto) y Spearman
cor(perfil_clientes$compras, perfil_clientes$gasto)
cor(perfil_clientes$compras, perfil_clientes$gasto, method = "spearman")
cor(perfil_clientes$edad, perfil_clientes$gasto)
\end{rcode}
\begin{rsalida}
[1] 0.7927073
[1] 0.7519843
[1] -0.07345163
\end{rsalida}

El número de compras tiene una correlación fuerte con el gasto (0.79 con
Pearson y 0.75 con Spearman): los clientes que más gastan son,
sobre todo, los que compran más veces. La edad, en cambio, casi no tiene
relación con el gasto ($-0.07$).

Para ver la diferencia entre ambos métodos, observa este ejemplo pequeño
creado a mano: diez tiendas donde las ventas suben siempre con el número de
visitantes, salvo que la décima tienda tuvo una venta corporativa enorme.

\begin{rcode}
visitantes <- c(100, 120, 150, 170, 200, 220, 250, 270, 300, 320)
ventas_dia <- c(10, 12, 14, 17, 19, 22, 24, 27, 29, 300)  # la última: atípica

cor(visitantes, ventas_dia)                        # Pearson
cor(visitantes, ventas_dia, method = "spearman")   # Spearman
\end{rcode}
\begin{rsalida}
[1] 0.5726456
[1] 1
\end{rsalida}

La relación es perfectamente monótona (más visitantes, siempre más ventas),
y Spearman lo detecta con $\rho = 1$. Pearson da 0.57 porque el dato
atípico rompe la línea recta. \textbf{Regla práctica}: usa Pearson cuando la
relación se vea lineal y no haya atípicos fuertes; usa Spearman con datos
asimétricos, con atípicos o con escalas ordinales (calificaciones de 1 a 5).

\subsection{Matriz de correlación}

Cuando tienes muchas variables, \codigo{cor()} aplicado a un data frame
numérico calcula todas las correlaciones a la vez: una
\termino{matriz de correlación}.

\begin{rcode}
matriz_cor <- perfil_clientes %>%
  select(compras, unidades, gasto, ticket, edad, antiguedad) %>%
  cor()

round(matriz_cor, 2)
\end{rcode}
\begin{rsalida}
           compras unidades gasto ticket  edad antiguedad
compras       1.00     0.91  0.79  -0.09 -0.02       0.01
unidades      0.91     1.00  0.87   0.14 -0.01       0.03
gasto         0.79     0.87  1.00   0.46 -0.07       0.05
ticket       -0.09     0.14  0.46   1.00 -0.14       0.05
edad         -0.02    -0.01 -0.07  -0.14  1.00       0.01
antiguedad    0.01     0.03  0.05   0.05  0.01       1.00
\end{rsalida}

La diagonal siempre vale 1 (cada variable consigo misma) y la matriz es
simétrica: la correlación de \codigo{compras} con \codigo{gasto} es la misma
que la de \codigo{gasto} con \codigo{compras}. Leer una tabla así es
cansado; \paquete{corrplot} la convierte en una imagen:

\begin{rcode}
corrplot(matriz_cor,
         method = "color",          # celdas coloreadas
         type = "upper",            # solo el triángulo superior
         addCoef.col = "black",     # escribe el coeficiente
         number.cex = 0.8,          # tamaño de los números
         tl.col = "black",          # color de los nombres
         tl.srt = 45,               # nombres inclinados
         col = colorRampPalette(c("#e34948", "white", "#2a78d6"))(200),
         diag = FALSE)              # sin la diagonal de unos
\end{rcode}

\begin{figure}[H]
\centering
\includegraphics[width=0.75\textwidth]{m04-corrplot}
\caption{Matriz de correlación del perfil de clientes. Escala divergente:
azul = correlación positiva, rojo = negativa, blanco = sin relación.}
\label{fig:m04-corrplot}
\end{figure}

La figura~\ref{fig:m04-corrplot} usa una paleta \emph{divergente} (rojo,
blanco, azul) porque la correlación tiene un centro natural en cero. Se lee
de un vistazo: el gasto se explica por la frecuencia (\codigo{compras},
0.79) y por las \codigo{unidades} (0.87); el \codigo{ticket} promedio tiene
una relación moderada con el gasto (0.46), y la edad y la antigüedad no se
relacionan con nada. Para la directora esto significa que \textbf{aumentar la
frecuencia de compra} (que el cliente regrese) pesa más que la edad o la
antigüedad del cliente.

\subsection{¿La correlación es real? \texorpdfstring{\codigo{cor.test()}}{cor.test()}}

Una correlación calculada con pocos datos puede salir distinta de cero por
pura casualidad. La función \codigo{cor.test()} responde: \emph{si en
realidad no hubiera ninguna relación, ¿qué tan raro sería observar una
correlación como esta?} (en la sección~\ref{sec:m4-hipotesis} verás esta
lógica con detalle).

\begin{rcode}
cor.test(perfil_clientes$edad, perfil_clientes$gasto)
\end{rcode}
\begin{rsalida}
	Pearson's product-moment correlation

data:  perfil_clientes$edad and perfil_clientes$gasto
t = -1.0311, df = 196, p-value = 0.3038
alternative hypothesis: true correlation is not equal to 0
95 percent confidence interval:
 -0.21073472  0.06667273
sample estimates:
        cor
-0.07345163
\end{rsalida}

Lo que importa de esta salida:
\begin{itemize}
  \item \codigo{cor = -0.073}: la correlación de la muestra.
  \item \codigo{p-value = 0.3038}: una probabilidad de 30\%. Si la edad no
        tuviera ninguna relación con el gasto, sería muy común ver una
        correlación como esta solo por azar. Conclusión: \textbf{no hay
        evidencia} de que la edad se relacione con el gasto.
  \item \codigo{95 percent confidence interval}: el valor real de la
        correlación podría estar entre $-0.21$ y $0.07$; el intervalo
        incluye el cero.
\end{itemize}

\begin{hazlotu}
Aplica \codigo{cor.test()} a \codigo{compras} y \codigo{gasto}. ¿Qué
p-valor obtienes? ¿El intervalo de confianza incluye el cero? Repite con
\codigo{method = "spearman"}.
\end{hazlotu}

\subsection{Correlación no implica causalidad}
\label{sec:m4-causalidad}

Que dos variables se muevan juntas \textbf{no significa} que una cause a la
otra. Puede haber una tercera variable, llamada \termino{variable de
confusión}, que mueve a las dos. Para verlo con claridad \textbf{simulamos}
datos donde conocemos la verdad: en una cadena de tiendas, la temperatura
del día aumenta tanto la venta de helados como la de ventiladores, pero los
helados \emph{no} influyen en los ventiladores.

\begin{rcode}
set.seed(123)
dias_verano <- tibble(
  temperatura = runif(90, min = 18, max = 38),           # grados
  helados     = 20 + 3 * temperatura + rnorm(90, sd = 8),
  ventiladores = 5 + 1.5 * temperatura + rnorm(90, sd = 5)
)

# Ventas de helados vs ventiladores: correlación alta...
cor(dias_verano$helados, dias_verano$ventiladores)
\end{rcode}
\begin{rsalida}
[1] 0.8160674
\end{rsalida}

Una correlación de 0.82 es muy fuerte. Un análisis ingenuo diría: ``pongamos
los ventiladores junto a los helados, porque vender helados impulsa la venta
de ventiladores''. Pero sabemos (porque así los simulamos) que la única causa
común es la temperatura. En la sección~\ref{sec:m4-regresion} verás cómo la
regresión múltiple permite ``controlar'' por la temperatura y comprobar que,
a igual temperatura, los helados no explican los ventiladores.

\begin{negocio}{Correlaciones engañosas en la empresa}
\begin{itemize}
  \item ``Las tiendas con más empleados venden más'': tal vez las tiendas
        grandes tienen más empleados \emph{y} más clientes.
  \item ``Los clientes que usan la app gastan más'': quizá los clientes que
        ya gastaban mucho son los que descargan la app.
  \item ``Cuando subimos la publicidad subieron las ventas'': ¿o se subió la
        publicidad justo en diciembre, cuando las ventas suben solas?
\end{itemize}
Para demostrar causalidad se necesitan experimentos controlados (como las
pruebas A/B de la sección~\ref{sec:m4-ab}) o técnicas que controlen las
variables de confusión.
\end{negocio}

% ----------------------------------------------------------------------------
\section{Probabilidad y distribuciones útiles}
\label{sec:m4-probabilidad}

Una \termino{distribución de probabilidad} es un modelo que describe qué
valores puede tomar algo incierto y con qué probabilidad. No necesitas
memorizar fórmulas: R trae funciones con una convención de nombres muy
sencilla (tabla~\ref{tab:m4-distribuciones}).

\begin{table}[H]
\centering
\small
\caption{Prefijos de las funciones de distribución en R.}
\label{tab:m4-distribuciones}
\begin{tabularx}{\textwidth}{@{}l X l@{}}
\toprule
\textbf{Prefijo} & \textbf{Qué calcula} & \textbf{Ejemplos} \\
\midrule
\codigo{d} & Probabilidad de un valor exacto (densidad) & \codigo{dbinom()},
  \codigo{dpois()} \\
\codigo{p} & Probabilidad acumulada: $P(X \leq x)$ & \codigo{pnorm()},
  \codigo{pbinom()} \\
\codigo{q} & El valor que deja por debajo una probabilidad dada (cuantil) &
  \codigo{qnorm()} \\
\codigo{r} & Números aleatorios que siguen esa distribución & \codigo{rnorm()},
  \codigo{rbinom()} \\
\bottomrule
\end{tabularx}
\end{table}

\subsection{Distribución normal: ventas semanales}

La \termino{distribución normal} es la famosa ``campana de Gauss'':
simétrica, con la mayoría de los datos cerca de la media y cada vez menos
hacia los extremos. Queda definida por dos números: la media ($\mu$) y la
desviación estándar ($\sigma$). Una regla útil: en una normal, cerca del
68\% de los datos cae a menos de una desviación estándar de la media y el
95\% a menos de dos.

El ticket individual no es normal (tiene cola a la derecha), pero las
\emph{sumas} de muchas ventas tienden a serlo. Esto se debe al
\termino{teorema del límite central}, uno de los resultados más importantes
de la estadística: al sumar o promediar muchos datos independientes, el
resultado se parece cada vez más a una normal, sin importar la forma de los
datos originales. Comprobémoslo con las ventas semanales:

\begin{rcode}
ventas_semanales <- transacciones %>%
  # semanas completas de lunes a domingo (del 2 de enero al 31 de dic.)
  filter(fecha >= as.Date("2023-01-02")) %>%
  mutate(semana = floor_date(fecha, unit = "week", week_start = 1)) %>%
  group_by(semana) %>%
  summarise(ventas = sum(total_transaccion), .groups = "drop")

media_sem <- mean(ventas_semanales$ventas)
sd_sem    <- sd(ventas_semanales$ventas)
nrow(ventas_semanales)
c(media = media_sem, desv_est = sd_sem)
asimetria(ventas_semanales$ventas)
\end{rcode}
\begin{rsalida}
[1] 52
    media  desv_est
18444.936  4591.285
[1] 0.2219801
\end{rsalida}

Tenemos 52 semanas completas con venta media de \$18,445 y desviación
estándar de \$4,591. La asimetría es pequeña (0.22): la forma es casi
simétrica. La
figura~\ref{fig:m04-normal-semanal} compara el histograma con la curva
normal teórica.

\begin{rcode}
p_normal <- ggplot(ventas_semanales, aes(x = ventas)) +
  geom_histogram(aes(y = after_stat(density)), binwidth = 2000,
                 fill = color_principal, color = "white", alpha = 0.85) +
  # curva normal con la media y desviación de los datos
  stat_function(fun = dnorm, args = list(mean = media_sem, sd = sd_sem),
                color = color_resalte, linewidth = 0.9) +
  geom_vline(xintercept = 25000, linetype = "dashed", color = "#52514e") +
  annotate("text", x = 25300, y = 9e-05, hjust = 0, size = 3.3,
           label = "¿Más de $25,000?", color = "#52514e") +
  scale_x_continuous(labels = dollar) +
  scale_y_continuous(labels = NULL) +
  labs(title = "Ventas semanales de 2023 y curva normal",
       subtitle = "52 semanas; la curva naranja es Normal(18,445; 4,591)",
       x = "Ventas de la semana", y = "Densidad") +
  tema_libro()

p_normal
\end{rcode}

\begin{figure}[H]
\centering
\includegraphics[width=0.9\textwidth]{m04-normal-semanal}
\caption{Las ventas semanales se ajustan bien a una distribución normal.}
\label{fig:m04-normal-semanal}
\end{figure}

Con el modelo normal podemos responder preguntas de planeación:

\begin{ejemplo}{Probabilidades con la normal}
\begin{enumerate}
  \item ¿Qué probabilidad hay de vender más de \$25,000 en una semana (lo
        que activaría un bono para el equipo)?
  \item ¿Cuánto inventario (en pesos) debe haber para cubrir la demanda del
        90\% de las semanas?
  \item ¿Qué probabilidad hay de una semana ``mala'', con menos de \$12,000?
\end{enumerate}
\end{ejemplo}

\begin{rcode}
# 1. P(ventas > 25,000) = 1 - P(ventas <= 25,000)
1 - pnorm(25000, mean = media_sem, sd = sd_sem)

# Comprobación con los datos reales: proporción de semanas > 25,000
mean(ventas_semanales$ventas > 25000)

# 2. Venta semanal que solo se supera el 10% de las veces
qnorm(0.90, mean = media_sem, sd = sd_sem)

# 3. P(ventas < 12,000)
pnorm(12000, mean = media_sem, sd = sd_sem)
\end{rcode}
\begin{rsalida}
[1] 0.0766864
[1] 0.07692308
[1] 24328.9
[1] 0.08019928
\end{rsalida}

Respuestas: (1) hay un 7.7\% de probabilidad de superar \$25,000 en una
semana, unas 4 semanas al año; los datos reales coinciden (7.7\%, es decir,
4 de 52 semanas). (2) Con un inventario equivalente a \$24,329 se cubre la
demanda del 90\% de las semanas. (3) Una semana por debajo de \$12,000 tiene
una probabilidad de 8.0\%.

\begin{advertencia}
Verifica siempre que la normal sea razonable antes de usarla (histograma,
asimetría cercana a cero). Si hubieras usado la normal con los
\emph{tickets individuales}, que tienen cola larga, habrías calculado
probabilidades negativas o muy equivocadas para los tickets bajos.
\end{advertencia}

\subsection{Distribución binomial: conversiones de una campaña}

La \termino{distribución binomial} cuenta cuántos ``éxitos'' hay en $n$
intentos independientes, cuando cada intento tiene la misma probabilidad
$p$ de éxito. Ejemplos: cuántos de 200 correos generan una compra, cuántos
de 50 clientes aceptan una oferta, cuántas de 1,000 piezas salen
defectuosas. El número esperado de éxitos es $n \times p$.

\begin{ejemplo}{Campaña de correo a clientes premium}
La empresa enviará un correo con una oferta a 200 clientes. Por campañas
anteriores, sabe que el 5\% de quienes reciben el correo compran. ¿Cuántas
compras debe esperar? ¿Qué probabilidad hay de conseguir exactamente 10?
¿Y de conseguir 15 o más, que es lo que el gerente prometió?
\end{ejemplo}

\begin{rcode}
n_correos  <- 200     # intentos
p_conv     <- 0.05    # probabilidad de conversión

n_correos * p_conv                            # conversiones esperadas
dbinom(10, size = n_correos, prob = p_conv)   # P(exactamente 10)
pbinom(5, size = n_correos, prob = p_conv)    # P(5 o menos)
1 - pbinom(14, size = n_correos, prob = p_conv)   # P(15 o más)
\end{rcode}
\begin{rsalida}
[1] 10
[1] 0.1283574
[1] 0.0623425
[1] 0.07813442
\end{rsalida}

Lo esperado son 10 compras, pero conseguir \emph{exactamente} 10 solo tiene
un 12.8\% de probabilidad: el azar mueve el resultado. Hay un 6.2\% de
probabilidad de un resultado decepcionante de 5 o menos, y apenas un 7.8\%
de llegar a 15 o más. El gerente prometió algo poco probable: para esperar
15 compras con un 5\% de conversión habría que enviar 300 correos.

\begin{rcode}
# ¿Cuántos correos hay que enviar para tener 90% de probabilidad
# de conseguir al menos 15 compras?
envios <- seq(200, 600, by = 10)
prob_15_o_mas <- 1 - pbinom(14, size = envios, prob = p_conv)
min(envios[prob_15_o_mas >= 0.90])
\end{rcode}
\begin{rsalida}
[1] 400
\end{rsalida}

Para prometer 15 compras con un 90\% de seguridad hay que enviar al menos 400
correos, el doble de lo planeado. Este tipo de cálculo evita compromisos
imposibles con la dirección.

\subsection{Distribución de Poisson: llegadas de clientes}

La \termino{distribución de Poisson} cuenta cuántos eventos ocurren en un
intervalo (de tiempo o de espacio) cuando ocurren al azar a un ritmo
promedio constante $\lambda$ (``lambda''). Ejemplos: clientes que llegan a
la caja por hora, llamadas al centro de atención por minuto, quejas por día.
Una propiedad característica: en una Poisson la media y la varianza son
iguales ($\lambda$).

Nuestros datos permiten comprobarlo. Contemos cuántas transacciones hubo
cada uno de los 365 días de 2023 (incluidos los días con cero ventas):

\begin{rcode}
tx_por_dia <- tibble(fecha = seq(as.Date("2023-01-01"),
                                 as.Date("2023-12-31"), by = "day")) %>%
  left_join(count(transacciones, fecha), by = "fecha") %>%
  mutate(n = replace_na(n, 0))              # días sin ventas = 0

lambda <- mean(tx_por_dia$n)
c(media = lambda, varianza = var(tx_por_dia$n))

# Frecuencia observada vs la que predice una Poisson(lambda)
tibble(transacciones_dia = 0:6) %>%
  mutate(
    dias_observados = map_int(transacciones_dia,
                              ~ sum(tx_por_dia$n == .x)),
    dias_poisson = round(365 * dpois(transacciones_dia, lambda), 1)
  )
\end{rcode}
\begin{rsalida}
   media varianza
2.739726 2.544709
# A tibble: 7 x 3
  transacciones_dia dias_observados dias_poisson
              <int>           <int>        <dbl>
1                 0              20         23.6
2                 1              69         64.6
3                 2              85         88.5
4                 3              78         80.8
5                 4              63         55.3
6                 5              32         30.3
7                 6              14         13.8
\end{rsalida}

La media (2.74) y la varianza (2.54) son muy parecidas, y los días
observados con 0, 1, 2\ldots{} transacciones se parecen mucho a los que
predice la Poisson. \codigo{map\_int()} (del paquete \paquete{purrr}, parte
del \paquete{tidyverse}) aplica el conteo a cada valor de
\codigo{transacciones\_dia}.

\begin{ejemplo}{Personal en cajas}
En la sucursal Centro llegan en promedio 12 clientes por hora a las cajas
en la tarde. Cada cajero atiende cómodamente a 8 clientes por hora. Con dos
cajeros (capacidad de 16), ¿qué probabilidad hay de que en una hora lleguen
más clientes de los que se pueden atender? ¿Y con tres cajeros
(capacidad 24)?
\end{ejemplo}

\begin{rcode}
lambda_hora <- 12
1 - ppois(16, lambda = lambda_hora)   # P(más de 16): 2 cajeros
1 - ppois(24, lambda = lambda_hora)   # P(más de 24): 3 cajeros
dpois(0, lambda = lambda_hora)        # P(ningún cliente en la hora)
\end{rcode}
\begin{rsalida}
[1] 0.101291
[1] 0.0006856332
[1] 6.144212e-06
\end{rsalida}

Con dos cajeros, en un 10.1\% de las horas de la tarde habrá más clientes
que capacidad (se formarán filas, aproximadamente 1 de cada 10 horas). Con
tres cajeros ese riesgo baja a 0.07\%. La última línea muestra que una hora
sin ningún cliente es prácticamente imposible (0.0006\%; el
\codigo{e-06} significa ``por $10^{-6}$''). La decisión ya no es ``intuición
contra intuición'': es comparar el costo de un tercer cajero contra el costo
de las filas en el 10\% de las horas.

\begin{hazlotu}
Si el ritmo sube a 15 clientes por hora en diciembre, ¿qué probabilidad hay
de superar la capacidad con dos cajeros? ¿Y con tres?
\end{hazlotu}

% ----------------------------------------------------------------------------
\section{Muestreo e intervalos de confianza}
\label{sec:m4-intervalos}

Casi nunca tienes todos los datos que te interesan. Los 1,000 tickets de
2023 son una \termino{muestra} de todas las compras que tus clientes harán;
una encuesta a 400 clientes es una muestra de toda la clientela (la
\termino{población}). La pregunta clave es: \emph{si tomara otra muestra,
¿cuánto cambiaría mi resultado?}

\subsection{Error estándar}

El \termino{error estándar} (EE) mide cuánto variaría la media si repitieras
el muestreo muchas veces. Se calcula así:
\[
  EE = \frac{s}{\sqrt{n}}
\]
Fíjate en el $\sqrt{n}$ del denominador: con más datos, el error estándar
baja, pero despacio. Para reducir el error a la mitad necesitas
\emph{cuatro} veces más datos.

\begin{rcode}
n_tickets <- length(ticket)
error_estandar <- sd(ticket) / sqrt(n_tickets)
error_estandar
\end{rcode}
\begin{rsalida}
[1] 24.0448
\end{rsalida}

No confundas los dos conceptos: la desviación estándar (\$760) dice cuánto
varían los tickets \emph{individuales}; el error estándar (\$24) dice cuánto
variaría el \emph{promedio} de 1,000 tickets.

\subsection{Intervalo de confianza para el ticket promedio}

Un \termino{intervalo de confianza} (IC) es un rango de valores plausibles
para el verdadero valor de la población. El IC del 95\% para una media se
construye como:
\[
  \bar{x} \;\pm\; t \times EE
\]
donde $t$ es un número que sale de la \termino{distribución t de Student}
(muy parecida a la normal) y, con muestras grandes, vale casi 2. Hagámoslo
``a mano'' y luego con \codigo{t.test()}:

\begin{rcode}
t_critico <- qt(0.975, df = n_tickets - 1)   # valor t para 95%
t_critico
media_ticket - t_critico * error_estandar    # límite inferior
media_ticket + t_critico * error_estandar    # límite superior

# Lo mismo en una línea:
t.test(ticket)$conf.int
\end{rcode}
\begin{rsalida}
[1] 1.962341
[1] 912.8086
[1] 1007.177
[1]  912.8086 1007.1768
attr(,"conf.level")
[1] 0.95
\end{rsalida}

\begin{concepto}{Cómo interpretar un intervalo de confianza}
Con los datos de 2023, el ticket promedio está entre \$912.81 y \$1,007.18
con una confianza del 95\%. La interpretación correcta es: \emph{si
repitiéramos este muestreo muchas veces y calculáramos un intervalo con cada
muestra, el 95\% de esos intervalos contendría el verdadero ticket
promedio}. En la práctica, dilo así: ``estimamos un ticket promedio de \$960,
con un margen de error de $\pm$\$47''.

\textbf{No es correcto} decir ``el 95\% de los tickets está entre \$913 y
\$1,007'': el intervalo habla del \emph{promedio}, no de los tickets
individuales (que van de \$23 a \$2,992).
\end{concepto}

\begin{negocio}{El margen de error en las decisiones}
Si el presupuesto de 2024 supone 12,000 transacciones, el ingreso esperado
es $12{,}000 \times 960 \approx$ 11.5 millones, pero con el IC podrías
presentar un rango: entre $12{,}000 \times 913 \approx$ 10.95 y
$12{,}000 \times 1{,}007 \approx$ 12.09 millones. Presentar un rango en
vez de un número único es más honesto y ayuda a planear escenarios.
\end{negocio}

\subsection{Bootstrap: intervalos sin fórmulas}

¿Y si quieres un intervalo para la \emph{mediana}, o para cualquier medida
que no tenga una fórmula sencilla? El \termino{bootstrap} es una técnica
ingeniosa: simula ``otras muestras'' tomando, con reemplazo, 1,000 tickets al
azar de tus propios 1,000 tickets. Algunos tickets saldrán repetidos y otros
no saldrán; así cada remuestra es ligeramente distinta. Repites esto miles de
veces, calculas la mediana de cada remuestra y te quedas con el 95\% central
de esos resultados.

\begin{rcode}
set.seed(123)   # para que el resultado sea reproducible
medianas_boot <- replicate(5000, {
  remuestra <- sample(ticket, size = length(ticket), replace = TRUE)
  median(remuestra)
})

length(medianas_boot)                           # 5,000 medianas
quantile(medianas_boot, probs = c(0.025, 0.975))   # IC 95% (percentiles)
\end{rcode}
\begin{rsalida}
[1] 5000
  2.5%  97.5%
641.23 796.43
\end{rsalida}

\codigo{replicate(5000, \{...\})} repite 5,000 veces el código entre llaves
y guarda cada resultado. El IC del 95\% para la mediana del ticket va de unos
\$641 a \$796. La figura~\ref{fig:m04-bootstrap} muestra las 5,000 medianas
simuladas.

\begin{rcode}
ic_boot <- quantile(medianas_boot, probs = c(0.025, 0.975))

p_boot <- ggplot(tibble(mediana = medianas_boot), aes(x = mediana)) +
  geom_histogram(binwidth = 5, fill = color_principal, color = "white") +
  geom_vline(xintercept = ic_boot, color = color_resalte,
             linewidth = 0.8, linetype = "dashed") +
  geom_vline(xintercept = mediana_ticket, color = "#0b0b0b",
             linewidth = 0.8) +
  scale_x_continuous(labels = dollar) +
  labs(title = "Distribución bootstrap de la mediana del ticket",
       subtitle = "5,000 remuestras; líneas naranjas = IC del 95%",
       x = "Mediana del ticket en cada remuestra",
       y = "Frecuencia") +
  tema_libro()

p_boot
\end{rcode}

\begin{figure}[H]
\centering
\includegraphics[width=0.9\textwidth]{m04-bootstrap}
\caption{Medianas de 5,000 remuestras bootstrap. La línea negra es la
mediana observada; entre las líneas naranjas queda el 95\% de las
remuestras.}
\label{fig:m04-bootstrap}
\end{figure}

\begin{hazlotu}
Usa el bootstrap para obtener un IC del 95\% del \emph{percentil 90} del
ticket (cambia \codigo{median(remuestra)} por
\codigo{quantile(remuestra, 0.9)}). ¿Es más ancho o más angosto que el de la
mediana? ¿Por qué crees que pasa?
\end{hazlotu}

% ----------------------------------------------------------------------------
\section{Pruebas de hipótesis para negocios}
\label{sec:m4-hipotesis}

Llegamos a la pregunta que aparece en todas las juntas: \emph{``¿esta
diferencia es real o es casualidad?''}. Los clientes premium gastan en
promedio \$977 por ticket y los demás \$951. ¿Es suficiente para decir que
los premium gastan más? Con otra muestra de tickets, la diferencia podría
desaparecer o invertirse. Las \termino{pruebas de hipótesis} dan un método
para decidir.

\subsection{La lógica en cinco conceptos}

\begin{concepto}{Hipótesis nula, alternativa y p-valor}
\begin{itemize}
  \item \termino{Hipótesis nula} ($H_0$): la postura ``aburrida'': no hay
        diferencia, no hay efecto. Ejemplo: ``el ticket promedio de premium y
        no premium es igual''.
  \item \termino{Hipótesis alternativa} ($H_1$): lo que quieres demostrar:
        ``el ticket promedio es distinto''.
  \item \termino{p-valor}: la probabilidad de observar una diferencia
        \emph{tan grande o más} que la de tus datos \textbf{si $H_0$ fuera
        cierta}. Un p-valor pequeño significa: ``sería muy raro ver esto por
        puro azar''.
  \item \termino{Nivel de significancia} ($\alpha$): el umbral que fijas
        \emph{antes} de ver los datos, casi siempre 0.05 (5\%). Si
        $p < \alpha$ rechazas $H_0$ y dices que el resultado es
        \termino{estadísticamente significativo}. Si $p \geq \alpha$,
        \emph{no} rechazas $H_0$: no hay evidencia suficiente.
\end{itemize}
\end{concepto}

Una analogía útil es un juicio: $H_0$ es ``el acusado es inocente''. Solo lo
declaras culpable si la evidencia es muy fuerte (p-valor pequeño). Si la
evidencia no alcanza, el veredicto es ``no culpable'', que \textbf{no es lo
mismo} que ``inocente demostrado''. Igual en estadística: no rechazar $H_0$
no demuestra que no haya diferencia; solo dice que tus datos no la
demuestran.

\begin{table}[H]
\centering
\small
\caption{Los dos errores posibles en una prueba de hipótesis.}
\label{tab:m4-errores}
\begin{tabularx}{\textwidth}{@{}l X X@{}}
\toprule
 & \textbf{En realidad $H_0$ es cierta} & \textbf{En realidad $H_0$ es falsa} \\
\midrule
\textbf{Rechazas $H_0$} & \termino{Error tipo I} (falso positivo): lanzas a
  todo el país una campaña que en realidad no funciona. Su probabilidad es
  $\alpha$. & Acierto: detectas un efecto real. \\
\textbf{No rechazas $H_0$} & Acierto: no te dejas engañar por el azar. &
  \termino{Error tipo II} (falso negativo): descartas una campaña que sí
  funcionaba. Se reduce con muestras más grandes. \\
\bottomrule
\end{tabularx}
\end{table}

\begin{advertencia}
\textbf{Significancia estadística no es lo mismo que importancia práctica.}
Con millones de datos, una diferencia de 50 centavos en el ticket puede ser
``significativa'' ($p < 0.05$) y aun así no importarle a nadie. Con pocos
datos, una diferencia enorme puede no ser significativa. Reporta siempre el
\emph{tamaño} de la diferencia (en pesos o en puntos porcentuales) junto con
el p-valor, y pregúntate: ¿esta diferencia cambia alguna decisión?
\end{advertencia}

\subsection{Prueba t: ¿los clientes premium gastan más?}

La \termino{prueba t de Student} compara las medias de dos grupos. Primero
unimos las transacciones con los clientes (Módulo~\ref{cap:manipulacion})
para saber quién es premium:

\begin{rcode}
tx_clientes <- transacciones %>%
  left_join(clientes, by = "cliente_id")

tx_clientes %>%
  group_by(es_premium) %>%
  summarise(tickets = n(),
            media   = mean(total_transaccion),
            mediana = median(total_transaccion),
            desv    = sd(total_transaccion),
            .groups = "drop")
\end{rcode}
\begin{rsalida}
# A tibble: 2 x 5
  es_premium tickets media mediana  desv
  <lgl>        <int> <dbl>   <dbl> <dbl>
1 FALSE          653  951.    672.  775.
2 TRUE           347  977.    803.  732.
\end{rsalida}

Planteamos: $H_0$: el ticket promedio es igual en ambos grupos;
$H_1$: es distinto. Nivel de significancia $\alpha = 0.05$.

\begin{rcode}
prueba_premium <- t.test(total_transaccion ~ es_premium,
                         data = tx_clientes)
prueba_premium
\end{rcode}
\begin{rsalida}
	Welch Two Sample t-test

data:  total_transaccion by es_premium
t = -0.52255, df = 741.31, p-value = 0.6014
alternative hypothesis: true difference in means between group FALSE and group TRUE is not equal to 0
95 percent confidence interval:
 -123.43528   71.53793
sample estimates:
mean in group FALSE  mean in group TRUE
           950.9885            976.9372
\end{rsalida}

Lectura línea por línea:
\begin{itemize}
  \item \codigo{Welch Two Sample t-test}: R usa por defecto la versión de
        Welch, que no supone que los dos grupos tengan la misma varianza. Es
        la opción segura.
  \item \codigo{t = -0.52255}: el \termino{estadístico t}, la diferencia de
        medias medida en errores estándar. Valores lejos de cero (más allá de
        $\pm 2$, aproximadamente) indican diferencias claras; $-0.52$ está
        muy cerca de cero.
  \item \codigo{p-value = 0.6014}: si premium y no premium gastaran lo mismo,
        veríamos una diferencia como esta el 60\% de las veces. Como
        $0.60 > 0.05$, \textbf{no rechazamos $H_0$}.
  \item \codigo{95 percent confidence interval}: la diferencia real
        (no premium menos premium) está entre $-123$ y $+72$ pesos. Incluye
        el cero.
  \item \codigo{sample estimates}: las medias de cada grupo.
\end{itemize}

Puedes extraer las piezas para usarlas en un reporte:

\begin{rcode}
prueba_premium$p.value
prueba_premium$conf.int
\end{rcode}
\begin{rsalida}
[1] 0.6014434
[1] -123.43528   71.53793
attr(,"conf.level")
[1] 0.95
\end{rsalida}

\begin{negocio}{``No hay evidencia'' es un resultado valioso}
Conclusión para la directora: \emph{con los datos de 2023 no hay evidencia
de que los clientes premium tengan un ticket promedio distinto al de los
demás; la diferencia observada (\$26) es compatible con el azar}. Esto no es
un fracaso del análisis: es información para decidir. Si la membresía premium
cuesta dinero (descuentos, envíos gratis), hay que preguntarse qué obtiene
la empresa a cambio. Quizá los premium compran con más \emph{frecuencia}
aunque no gasten más por ticket: esa es la siguiente hipótesis a revisar
(ejercicio~\ref{ej:m4-6}).
\end{negocio}

\subsection{Prueba A/B de proporciones: campaña A vs campaña B}
\label{sec:m4-ab}

Una \termino{prueba A/B} es un experimento: divides al azar a tus clientes en
dos grupos, a cada uno le muestras una versión distinta (de un correo, una
página, un precio) y comparas los resultados. Como la asignación es al azar,
las dos poblaciones son iguales en todo lo demás, y la diferencia se puede
atribuir a la versión. Es la forma más limpia de demostrar
\emph{causalidad} en un negocio.

\begin{ejemplo}{¿Qué asunto de correo convierte más?}
Mercadotecnia envió 2,000 correos con el asunto A (``Ofertas de la
semana'') y 2,000 con el asunto B (``Tu cupón personal te espera''). Como
queremos saber si el método funciona, \textbf{simulamos} los resultados con
tasas de conversión conocidas: 5\% para A y 6.5\% para B. En la vida real no
conocerías estas tasas: justamente eso es lo que la prueba intenta
averiguar.
\end{ejemplo}

\begin{rcode}
set.seed(123)
campana <- tibble(
  version  = rep(c("A", "B"), each = 2000),
  # 1 = compró, 0 = no compró; tasas reales: A 5%, B 6.5%
  convirtio = c(rbinom(2000, size = 1, prob = 0.050),
                rbinom(2000, size = 1, prob = 0.065))
)

resumen_ab <- campana %>%
  group_by(version) %>%
  summarise(enviados    = n(),
            conversiones = sum(convirtio),
            tasa        = mean(convirtio),
            .groups = "drop")
resumen_ab
\end{rcode}
\begin{rsalida}
# A tibble: 2 x 4
  version enviados conversiones   tasa
  <chr>      <int>        <int>  <dbl>
1 A           2000          101 0.0505
2 B           2000          130 0.065
\end{rsalida}

$H_0$: las dos versiones tienen la misma tasa de conversión. $H_1$: las tasas
son distintas. La función \codigo{prop.test()} recibe el número de éxitos y
el número de intentos de cada grupo:

\begin{rcode}
prueba_ab <- prop.test(x = resumen_ab$conversiones,
                       n = resumen_ab$enviados)
prueba_ab
\end{rcode}
\begin{rsalida}
	2-sample test for equality of proportions with continuity
	correction

data:  resumen_ab$conversiones out of resumen_ab$enviados
X-squared = 3.602, df = 1, p-value = 0.05771
alternative hypothesis: two.sided
95 percent confidence interval:
 -0.0294509801  0.0004509801
sample estimates:
prop 1 prop 2
0.0505 0.0650
\end{rsalida}

Lectura: la versión A convirtió 5.05\% y la B 6.50\%, pero el p-valor es
0.058, apenas \emph{mayor} a 0.05. Formalmente \textbf{no rechazamos $H_0$}:
con 2,000 correos por grupo, una diferencia así todavía podría deberse al
azar. El intervalo de confianza de la diferencia (A menos B) va de $-2.9$ a
$+0.05$ puntos porcentuales e incluye el cero.

Aquí hay una lección importante. Nosotros \emph{sabemos} que B es mejor
(así lo simulamos), y aun así la prueba no lo detectó: acabamos de cometer
un \textbf{error tipo II}. ¿La causa? Una muestra demasiado pequeña para
detectar una diferencia de 1.5 puntos. La función \codigo{power.prop.test()}
calcula cuántos envíos se necesitan para detectar una diferencia con una
probabilidad dada (la \termino{potencia}, típicamente 80\%):

\begin{rcode}
power.prop.test(p1 = 0.05, p2 = 0.065, power = 0.80)$n
\end{rcode}
\begin{rsalida}
[1] 3779.801
\end{rsalida}

Se necesitan unos 3,780 correos \emph{por grupo}. Repetimos el experimento
(de nuevo simulado, con las mismas tasas reales) con 4,000 correos por
versión:

\begin{rcode}
set.seed(123)
campana_4000 <- tibble(
  version   = rep(c("A", "B"), each = 4000),
  convirtio = c(rbinom(4000, size = 1, prob = 0.050),
                rbinom(4000, size = 1, prob = 0.065))
)

resumen_4000 <- campana_4000 %>%
  group_by(version) %>%
  summarise(enviados = n(), conversiones = sum(convirtio),
            .groups = "drop")

prop.test(x = resumen_4000$conversiones, n = resumen_4000$enviados)
\end{rcode}
\begin{rsalida}
	2-sample test for equality of proportions with continuity
	correction

data:  resumen_4000$conversiones out of resumen_4000$enviados
X-squared = 5.0351, df = 1, p-value = 0.02484
alternative hypothesis: two.sided
95 percent confidence interval:
 -0.022041478 -0.001458522
sample estimates:
 prop 1  prop 2
0.04975 0.06150
\end{rsalida}

Ahora $p = 0.025 < 0.05$: rechazamos $H_0$ y concluimos que el asunto B
convierte más (6.15\% contra 4.98\%). El intervalo (que R calcula como A menos B,
por eso sale negativo) dice que B supera a A por entre 0.15 y 2.2 puntos
porcentuales. Y la \emph{significancia práctica} es
clara: un punto más de conversión sobre 100,000 correos son unas 1,000
compras adicionales.

\begin{consejo}
Decide el tamaño de la prueba A/B \emph{antes} de lanzarla (con
\codigo{power.prop.test()}) y no la detengas en cuanto el p-valor baje de
0.05: si revisas el resultado cada hora y paras cuando ``sale'', aumentas
mucho la probabilidad de un falso positivo (error tipo I).
\end{consejo}

\subsection{ANOVA: ¿el ticket cambia entre categorías?}

Cuando quieres comparar las medias de \emph{más de dos} grupos, no hagas
muchas pruebas t de dos en dos (con 10 categorías serían 45 pruebas, y
alguna saldría ``significativa'' por azar). Usa el \termino{análisis de
varianza} (ANOVA). $H_0$: todas las medias son iguales; $H_1$: al menos una
es distinta. La idea es comparar cuánto varían las medias \emph{entre}
grupos con cuánto varían los datos \emph{dentro} de cada grupo.

\begin{rcode}
tx_productos <- transacciones %>%
  left_join(productos, by = "producto_id")

ticket_categoria <- tx_productos %>%
  group_by(categoria) %>%
  summarise(n     = n(),
            media = mean(total_transaccion),
            ee    = sd(total_transaccion) / sqrt(n()),
            .groups = "drop") %>%
  arrange(desc(media))
ticket_categoria
\end{rcode}
\begin{rsalida}
# A tibble: 10 x 4
   categoria       n media    ee
   <chr>       <int> <dbl> <dbl>
 1 Juguetes      118 1060.  77.1
 2 Alimentos     115 1057.  69.9
 3 Automotriz    105  991.  76.5
 4 Hogar          64  991.  90.7
 5 Belleza       187  974.  58.8
 6 Electrónica   137  968.  61.7
 7 Deportes       92  937.  78.3
 8 Libros         70  817.  85.2
 9 Ropa           38  791. 127.
10 Mascotas       74  779.  70.5
\end{rsalida}

Los tickets promedio van de \$779 (Mascotas) a \$1,060 (Juguetes). ¿Es una
diferencia real? Veamos las medias con su intervalo de confianza
aproximado (media $\pm$ 2 errores estándar):

\begin{rcode}
p_anova <- ggplot(ticket_categoria,
                  aes(x = media, y = reorder(categoria, media))) +
  geom_vline(xintercept = mean(ticket), linetype = "dashed",
             color = color_gris) +
  geom_errorbarh(aes(xmin = media - 2 * ee, xmax = media + 2 * ee),
                 height = 0.25, color = color_principal,
                 linewidth = 0.8) +
  geom_point(size = 2.5, color = color_principal) +
  scale_x_continuous(labels = dollar) +
  labs(title = "Ticket promedio por categoría con IC del 95%",
       subtitle = "Los intervalos se enciman mucho; línea gris = promedio general",
       x = "Ticket promedio", y = NULL) +
  tema_libro()

p_anova
\end{rcode}

\begin{figure}[H]
\centering
\includegraphics[width=0.9\textwidth]{m04-anova-categorias}
\caption{Ticket promedio por categoría. Cuando los intervalos se enciman
tanto, es difícil afirmar que alguna categoría sea distinta.}
\label{fig:m04-anova-categorias}
\end{figure}

\begin{rcode}
modelo_anova <- aov(total_transaccion ~ categoria, data = tx_productos)
summary(modelo_anova)
\end{rcode}
\begin{rsalida}
             Df    Sum Sq Mean Sq F value Pr(>F)
categoria     9   7466897  829655   1.441  0.166
Residuals   990 570107428  575866
\end{rsalida}

La columna clave es \codigo{Pr(>F)}, el p-valor: 0.166. Como es mayor a
0.05, \textbf{no hay evidencia} de que el ticket promedio dependa de la
categoría. El \codigo{F value} (1.441) compara la variación entre categorías
con la variación dentro de ellas; valores cercanos a 1 indican que las
categorías no se distinguen más de lo que se distinguen los tickets entre
sí. Hallazgo para la directora: las categorías difieren en \emph{volumen}
(Belleza tiene 187 tickets y Ropa 38), no en ticket promedio.

\subsubsection{Cuando el ANOVA sí es significativo: TukeyHSD}

Probemos ahora con el segmento de edad del cliente (Joven, Adulto, Maduro,
Senior):

\begin{rcode}
modelo_segmento <- aov(total_transaccion ~ segmento, data = tx_clientes)
summary(modelo_segmento)
\end{rcode}
\begin{rsalida}
             Df    Sum Sq Mean Sq F value Pr(>F)
segmento      3   4672766 1557589   2.708 0.0441 *
Residuals   996 572901558  575202
---
...
\end{rsalida}

Ahora $p = 0.044 < 0.05$: al menos un segmento tiene un ticket promedio
distinto. El asterisco junto al p-valor es la forma en que R marca la
significancia (una estrella: $p < 0.05$; dos: $p < 0.01$; tres:
$p < 0.001$); la leyenda completa aparece en la línea
\codigo{Signif. codes}, que aquí recortamos. ¿Cuál? El ANOVA no lo dice. Para eso se usa la prueba de
\termino{Tukey HSD} (\emph{Honest Significant Difference}), que compara
todos los pares de grupos ajustando los p-valores por hacer varias
comparaciones:

\begin{rcode}
TukeyHSD(modelo_segmento)
\end{rcode}
\begin{rsalida}
  Tukey multiple comparisons of means
    95% family-wise confidence level

Fit: aov(formula = total_transaccion ~ segmento, data = tx_clientes)

$segmento
                    diff        lwr       upr     p adj
Joven-Adulto    62.47256 -170.39680 295.34191 0.9007953
Maduro-Adulto -138.85507 -299.46454  21.75439 0.1172294
Senior-Adulto  -47.79558 -216.73288 121.14173 0.8858662
Maduro-Joven  -201.32763 -423.37427  20.71901 0.0912671
Senior-Joven  -110.26813 -338.41089 117.87463 0.5990701
Senior-Maduro   91.05950  -62.61673 244.73573 0.4229310
\end{rsalida}

Cada fila es una comparación: \codigo{diff} es la diferencia de medias,
\codigo{lwr} y \codigo{upr} su intervalo de confianza y \codigo{p adj} el
p-valor ajustado. Ningún par individual queda por debajo de 0.05; el más
cercano es Maduro contra Joven ($p = 0.09$): los clientes maduros gastan unos
\$201 menos por ticket que los jóvenes. Esto pasa a veces: el ANOVA detecta
que ``algo'' difiere, pero la evidencia es débil y no alcanza para señalar
un par concreto. \textbf{Conclusión honesta}: hay indicios de que el ticket
varía con la edad (los maduros gastan algo menos), pero la evidencia es
débil; conviene confirmarlo con los datos de 2024 antes de diseñar una
campaña por edad.

\subsection{Chi-cuadrada: ¿el método de pago depende de la tienda?}

Las pruebas anteriores comparan \emph{promedios} de una variable numérica.
Cuando las dos variables son \emph{categóricas}, se usa la \termino{prueba
chi-cuadrada de independencia} ($\chi^2$). $H_0$: las dos variables son
independientes (la mezcla de métodos de pago es la misma en todas las
tiendas); $H_1$: están relacionadas.

\begin{ejemplo}{Terminales de pago por sucursal}
Finanzas quiere saber si debe negociar terminales bancarias distintas para
cada sucursal o si todas las tiendas tienen la misma mezcla de métodos de
pago.
\end{ejemplo}

\begin{rcode}
tabla_pago <- table(transacciones$metodo_pago, transacciones$tienda_id)
tabla_pago
\end{rcode}
\begin{rsalida}
                   1  2  3  4  5  6  7  8  9 10
  Efectivo        26 36 35 29 33 43 28 31 27 27
  Tarjeta Crédito 37 37 35 29 38 28 32 29 45 40
  Tarjeta Débito  14 19 28 25 15 30 16 33 25 28
  Transferencia   10  7 11 11  9  8  9 13 11 13
\end{rsalida}

\begin{rcode}
prueba_chi <- chisq.test(tabla_pago)
prueba_chi
round(prueba_chi$expected[, 1:5], 1)   # conteos esperados (5 tiendas)
\end{rcode}
\begin{rsalida}
	Pearson's Chi-squared test

data:  tabla_pago
X-squared = 28.953, df = 27, p-value = 0.3632


                     1    2    3    4    5
  Efectivo        27.4 31.2 34.3 29.6 29.9
  Tarjeta Crédito 30.4 34.6 38.1 32.9 33.2
  Tarjeta Débito  20.3 23.1 25.4 21.9 22.1
  Transferencia    8.9 10.1 11.1  9.6  9.7
\end{rsalida}

El p-valor es 0.36: \textbf{no hay evidencia} de que el método de pago
dependa de la tienda. Las diferencias que ves en la tabla (por ejemplo, 13
transferencias en la tienda 8 contra 7 en la tienda 2) son las esperables
por azar. Los \codigo{expected} son los conteos que habría si las variables
fueran totalmente independientes; la prueba compara los observados con
estos. Decisión práctica: una sola negociación con el banco, con la misma
configuración de terminales para todas las sucursales.

\begin{advertencia}
La chi-cuadrada necesita conteos esperados de al menos 5 en casi todas las
celdas. Si hay muchas celdas con esperados pequeños, R muestra la advertencia
\codigo{Chi-squared approximation may be incorrect}. En ese caso agrupa
categorías poco frecuentes o usa \codigo{fisher.test()}.
\end{advertencia}

\subsection{¿Qué prueba uso?}

\begin{table}[H]
\centering
\small
\caption{Guía rápida para elegir una prueba de hipótesis.}
\label{tab:m4-que-prueba}
\begin{tabularx}{\textwidth}{@{}X X l@{}}
\toprule
\textbf{Pregunta de negocio} & \textbf{Tipo de datos} & \textbf{Función} \\
\midrule
¿El ticket promedio es distinto de \$900? & Una variable numérica vs un
  valor & \codigo{t.test(x, mu = 900)} \\
¿Premium y no premium gastan distinto? & Numérica + categórica de 2 grupos
  & \codigo{t.test(y \textasciitilde{} g)} \\
¿Mejoró la venta de cada tienda después del cambio? & Numérica, mismos
  sujetos antes y después & \codigo{t.test(antes, despues, paired = TRUE)} \\
¿Qué versión de la campaña convierte más? & Proporciones de 2 o más grupos &
  \codigo{prop.test()} \\
¿El ticket cambia entre categorías o tiendas? & Numérica + categórica de 3
  o más grupos & \codigo{aov()} + \codigo{TukeyHSD()} \\
¿El método de pago depende de la tienda? & Dos variables categóricas &
  \codigo{chisq.test()} \\
¿El gasto se relaciona con la edad? & Dos numéricas & \codigo{cor.test()} \\
Datos muy asimétricos o pocos datos (alternativa sin supuesto de
  normalidad) & Numérica + 2 grupos / 3+ grupos &
  \codigo{wilcox.test()} / \codigo{kruskal.test()} \\
\bottomrule
\end{tabularx}
\end{table}

\begin{hazlotu}
Usa \codigo{t.test()} para comparar el ticket de clientes hombres y mujeres
(columna \codigo{genero} de \codigo{tx\_clientes}). ¿El resultado es
significativo? ¿Cuál es la diferencia en pesos? Guarda tu respuesta: la
retomaremos en la sección de errores comunes.
\end{hazlotu}
